\begin{derivation}[局域标量顶角的配对计数]

正文\eqnrefs{eq:phi4-si-coupling-r10} 中显式写出的 $1/4!$ 不是图规则的附加约定，而是为了抵消同一个局域顶角上四个相同标量场的等价 Wick 配对数。逐个计数这 $4!$ 种收缩并截去外腿后，就得到正文\eqnrefs{eq:si-reduced-vertices-r10} 中的 $-\ii g_4$。



SI 相互作用密度已经固定为
\begin{equation*}
 \Lag_I(z)=-\frac{g_4}{4!\,\hbar c}\varphi^4(z).
\end{equation*}
因此 Dyson 一阶对四点 Green 函数的贡献与
\begin{equation*}
 \frac{1}{4!}
 \int\dd^4 z\,
 \langle0|\Tord\{
 \varphi(x_1)\varphi(x_2)\varphi(x_3)\varphi(x_4)
 \varphi^4(z)
 \}|0\rangle
 \end{equation*}
成正比；所有 $\hbar,c$ 的共同归一化已在
\eqnrefs{eq:si-reduced-vertices-r10}
的约化过程中统一提出。要得到连通树图，每个外部场必须与 $z$ 点四个场中的一个收缩。第一条外部场有 4 个选择，第二条剩 3 个，第三条剩 2 个，最后剩 1 个，所以同一个收缩乘积
\begin{equation*}
 \Delta_{1z}\Delta_{2z}\Delta_{3z}\Delta_{4z}
\end{equation*}
出现
\begin{equation*}
 4\times3\times2\times1=4!
\end{equation*}
次，恰好抵消拉氏量中的 $1/4!$。因此截肢外腿并提出统一国际单位制归一化后，约化顶角只剩
\begin{equation*}
 {-\ii g_4.}
\end{equation*}
$4!$ 不是被经验地“删掉”，而是被等价 Wick 收缩精确抵消。
\end{derivation}

\begin{derivation}[四点泡图的对称因子]
考虑 $\varphi^4$ 理论四点函数的二阶 Dyson 项
\begin{equation*}
 \frac{1}{2!}
 \left(-\frac{\ii g_4}{4!\,\hbar^2c^2}\right)^2
 \int\dd^4z\dd^4w\,
 \langle0|\mathcal T\{
 \varphi_1\varphi_2\varphi_3\varphi_4
 \varphi^4(z)\varphi^4(w)\}|0\rangle .
\end{equation*}
固定一个 $s$ 道拓扑：$x_1,x_2$ 接到一个顶角，$x_3,x_4$ 接到另一个顶角，两个顶角之间有两条内部线。计算其系数时，外腿分配和内部配对的数目要与 Dyson 展开的 $2!$、每个顶角的 $4!$ 一起核对。

先决定哪一个带积分标签的顶角接 $(x_1,x_2)$：有 $2$ 种选择。固定以后，在该顶角的四个同点场中，为有标签的 $x_1$ 选收缩对象有 $4$ 种，为 $x_2$ 再选有 $3$ 种，所以有 $4\times3=12$ 种。另一顶角给 $(x_3,x_4)$ 同样有 $12$ 种。每个顶角此时都剩两个场，两条跨顶角内部线的配对还有 $2!$ 种。因此产生同一个带标签 $s$ 道图的 Wick 收缩总数为
\begin{equation*}
 N_s=2\times12\times12\times2=576.
\end{equation*}
而 Dyson 与两个相互作用顶角的组合分母是
\begin{equation*}
 2!(4!)^2=2\times24^2=1152.
\end{equation*}
所以净组合系数严格为
\begin{equation*}
 \frac{N_s}{2!(4!)^2}=\frac{576}{1152}=\frac12.
\end{equation*}
这就是所谓泡图对称因子。它不是在画完图以后经验性补上的数，而是完整 Dyson--Wick 展开中“产生同一带标签图的收缩数”除以“指数展开和拉氏量中的排列分母”所得的商。

截肢四条外腿后，$s$ 道一圈振幅因而为
\begin{equation*}
 \ii\mathcal M_s^{(1)}
 =\frac12(-\ii g_4)^2
 \int_k^{(4)}
 \frac{\ii}{k^2-m^2c^2+\ii0^+}
 \frac{\ii}{(P-k)^2-m^2c^2+\ii0^+},
 \qquad P=p_1+p_2,
\end{equation*}
即正文\eqnrefs{eq:phi4-bubble-symmetry-factor-dp68}。$t$、$u$ 道来自外部标签的另外两种不等价分组，而不是这个 $1/2$ 内部对称因子的额外来源。
\end{derivation}

\begin{derivation}[局域电磁顶角中的动量守恒]
只看一次相互作用插入，并从三个自由场各取一个确定的平面波 Fourier 分量。把与 $x$ 无关的展开系数分别记为 $C_p,C_{p'},C_q$，则可严格写成
\begin{align*}
 \psi_p(x)&=C_pu(p)\ee^{-\ii p\cdot x/\hbar},\\
 \bar\psi_{p'}(x)&=C_{p'}^*\bar u(p')\ee^{+\ii p'\cdot x/\hbar},\\
 A_{q\mu}(x)&=C_q\varepsilon_\mu(q)\ee^{-\ii q\cdot x/\hbar},
\end{align*}
其中光子四动量 $q$ 取为流入顶角。三者代入\eqnrefs{eq:qed-S0-SI-dp7} 后，与 $x$ 有关的因子唯一合并成

\begin{equation*}
 \exp\!\left[
 \frac{\ii}{\hbar}(p'-p-q)\cdot x
 \right].
\end{equation*}
因此顶角时空积分是
\begin{align*}
 \int\dd^4 x\,
 \ee^{\frac{\ii}{\hbar}(p'-p-q)\cdot x}
 &=(2\pi\hbar)^4
 \delta^{(4)}(p'-p-q),\\
 &\Longrightarrow\qquad p+q=p'.
 \end{align*}
第一行只是全书 Fourier 约定下的 $\delta$-函数恒等式；第二行是它的支撑条件。也就是说，四动量守恒直接来自相互作用在时空中的平移不变性与局域积分。

内部线的四动量可以不满足 $q^2=m^2c^2$，因为它不是 LSZ 外态；但它连接的每个顶角仍各自携带一个这样的 \(\delta\) 函数。所以 ``虚粒子可以暂时违反能量守恒'' 并不是 QED 的数学陈述。真正被放宽的是内部动量的在壳约束，而不是顶角守恒律。
\end{derivation}

\begin{derivation}[高斯平均中的传播子与顶角规则]
先定义由自由 Dirac 与规范固定 Maxwell 二次作用量产生的归一化高斯平均

\begin{equation*}
 \langle\mathcal O\rangle_0
 \equiv
 \frac{\displaystyle
 \int\mathcal D\bar\psi\,\mathcal D\psi\,\mathcal DA\,
 \mathcal O\,\ee^{\ii S_0/\hbar}}
 {\displaystyle
 \int\mathcal D\bar\psi\,\mathcal D\psi\,\mathcal DA\,
 \ee^{\ii S_0/\hbar}}.
 \end{equation*}
自由高斯已在前面分别给出 Dirac 传播子 $S_F$ 与光子传播子 $D_F^{\mu\nu}$。相互作用理论的时序关联函数可以写成
\begin{equation*}
 {
 \langle\Tord\mathcal O\rangle
 =\frac{
 \left\langle\Tord\left\{\mathcal O
 \exp\!\left(\ii S_I/\hbar\right)\right\}\right\rangle_0}
 {\left\langle\Tord\exp\!\left(\ii S_I/\hbar\right)\right\rangle_0}.}
 \end{equation*}
分母的作用是消去与外部插入完全不连接的真空 bubbles；这正是连通 correlation 函数与真空归一化的关系。

为了看见基本顶角，不直接背诵 ``$-\ii e\gamma^\mu$''，而计算一个带两条费米子外插入和一条光子外插入的三点函数：
\begin{equation*}
 G_3^\mu(x,z;y)
 \equiv
 \left\langle
 \Tord\{\psi(x)A^\mu(z)\bar\psi(y)\}
 \right\rangle .
 \end{equation*}
把\eqnrefs{eq:qed-gell-mann-low-dp7} 的分子展开到 $q_{\mathrm e}$ 一阶，连通部分为
\begin{align*}
 G_{3,c}^{\mu,(1)}(x,z;y)
 ={}&-\frac{\ii q_{\mathrm e}}{\hbar}
 \int\dd^4 w\,
 \Big\langle\Tord\big\{
 \psi(x)A^\mu(z)\bar\psi(y)
 \bar\psi(w)\gamma^\nu\psi(w)A_\nu(w)
 \big\}\Big\rangle_{0,c}.
 \end{align*}
自由高斯中要让全部外部插入与相互作用点 $w$ 连通，只有一类收缩：
\begin{align*}
 \bigl[\psi(x)\,\bar\psi(w)\bigr]_{\rm contr}&\longrightarrow S_F(x-w),\\
 \bigl[\psi(w)\,\bar\psi(y)\bigr]_{\rm contr}&\longrightarrow S_F(w-y),\\
 \bigl[A^\mu(z)A_\nu(w)\bigr]_{\rm contr}&\longrightarrow D_F^{\mu}{}_{\nu}(z-w),
\end{align*}
下标 $\rm contr$ 表示选取这两个自由场的 Wick 收缩。
因此
\begin{equation*}
 {
 G_{3,c}^{\mu,(1)}(x,z;y)
 =-\frac{\ii q_{\mathrm e}}{\hbar}
 \int\dd^4 w\,
 S_F(x-w)\gamma^\nu S_F(w-y)
 D_F^{\mu}{}_{\nu}(z-w).}
 \end{equation*}
这条式子已经把图的三个组成部分全部暴露出来：$w$ 是顶角时空积分；三条自由收缩是三条传播子；唯一不属于自由传播的矩阵结构是 $\gamma^\nu$，它来自相互作用密度中的 Dirac 流。

Fourier 变换后，$w$ 积分产生四动量守恒 \(\delta\) 函数。再由 LSZ 去掉三条外部传播子极点，物理 SI 顶角因子就是
\begin{equation*}
 -\frac{\ii q_{\mathrm e}}{\hbar}\gamma^\nu,
 \end{equation*}
而把统一的 SI 场归一化常数提出后，约化规则变成前面定义的
\begin{equation*}
 {-\ii g_e\gamma^\nu,\qquad
 g_e=\frac{q_{\mathrm e}}{\sqrt{\varepsilon_0\hbar c}}.}
 \end{equation*}
所以 QED 顶角不是图论公设，而是 Maxwell--Dirac 作用量的一次微扰插入经过 Wick 收缩与 LSZ 后留下的截肢局域核。
\end{derivation}

\begin{derivation}[软光子因子的角分布]
先把\eqnrefs{eq:weinberg-soft-photon-factor-dp10} 写成速度形式。对第 $i$ 条外腿，

\begin{equation*}
 p_i^\mu=\frac{E_i}{c}(1,\boldsymbol\beta_i),
 \qquad
 k^\mu=\frac{\hbar\omega}{c}(1,\mathbf n),
 \qquad
 |\mathbf n|=1,
\end{equation*}
所以
\begin{equation*}
 p_i\!\cdot k
 =\frac{E_\ii\hbar\omega}{c^2}
 (1-\mathbf n\!\cdot\boldsymbol\beta_i).
\end{equation*}
对物理横向偏振求和时，只有速度相对于 $\mathbf n$ 的横向分量
\begin{equation*}
 \boldsymbol\beta_i^\perp
 =\boldsymbol\beta_i-\mathbf n(\mathbf n\!\cdot\boldsymbol\beta_i)
\end{equation*}
保留下来。使用采用上述协变态外线归一化的一光子相空间因子，并代入 $g_{\rm em}^2=4\pi\alpha$，真实软辐射相对于硬过程的领先概率为
\begin{equation*}
 \dd\mathcal R
 =\frac{\alpha}{4\pi^2}
 \frac{\dd\omega}{\omega}\,\dd\Omega
 \left|
 \sum_i\sigma_i^{\rm io}Q_i
 \frac{\boldsymbol\beta_i^\perp}
 {1-\mathbf n\!\cdot\boldsymbol\beta_i}
 \right|^2.
\end{equation*}
该式的量纲检查是直接的：$\alpha$、$\dd\omega/\omega$、$\dd\Omega$ 与速度比均无量纲。对
$\lambda_{\rm IR}/\hbar<\omega<\Delta E_{\rm det}/\hbar$ 积分即得\eqnrefs{eq:soft-radiator-definition-dp10}，其中径向积分为
\begin{equation*}
 \int_{\lambda_{\rm IR}/\hbar}^{\Delta E_{\rm det}/\hbar}
 \frac{\dd\omega}{\omega}
 =\ln\frac{\Delta E_{\rm det}}{\lambda_{\rm IR}}.
\end{equation*}
\end{derivation}

\begin{derivation}[虚软光子修正与实辐射的关系]
虚软光子连接外腿 $i,j$ 时，相邻费米子传播子的软极限为

\begin{equation*}
 \frac{\slashed p_i+\sigma_i^{\rm io}\slashed k+m_ic}
 {(p_i+\sigma_i^{\rm io} k)^2-m_i^2c^2+\ii0^+}
 =
 \frac{\slashed p_i+m_ic+\order(k)}
 {2\sigma_i^{\rm io} p_i\!\cdot k+\ii0^+}.
\end{equation*}
利用外腿 Dirac 方程把顶角两侧的分子约化为 $2p_i^\mu$，虚软修正的红外部分写成
\begin{equation*}
 \delta_{\rm virt}^{\rm soft}\mathcal M_0
 =-\frac12\mathcal M_0
 \sum_{ij}g_{\rm em}^2\sigma_i^{\rm io}\sigma_j^{\rm io}Q_iQ_j
 \int_{\rm soft}\frac{\dd^4 k}{(2\pi)^4}
 \frac{p_i\!\cdot p_j}
 {(k^2+\ii0^+)(p_i\!\cdot k+\ii0_i)(-p_j\!\cdot k+\ii0_j)}
 +\text{红外有限项}.
\end{equation*}
这里 $\ii0_i$ 记录外腿属于入射还是出射，决定相应的轮廓边界值。

对虚光子传播子取不连续度，
\begin{equation*}
 \frac1{k^2+\ii0^+}
 -\frac1{k^2-\ii0^+}
 =-2\pi\ii\,\delta(k^2).
\end{equation*}
正能割线进一步取 $\delta_+(k^2)=\Theta(k^0)\delta(k^2)$，于是四维圈积分中的虚光子线变成真实一光子相空间，并保留与 $\dd\mathcal R$ 相同的 $p_i,p_j,Q_iQ_j$ 成对核。光学定理把这部分解释为从无真实软光子通道流向真实软光子通道的概率，因此无辐射概率获得相反号。由此得到\eqnrefs{eq:virtual-real-ir-relation-dp10}。
\end{derivation}

\begin{derivation}[软光子实虚抵消与有限分辨率]
先把正文软辐射积分中的角变量完全积分掉，定义只依赖硬外腿运动学的非负系数
\begin{equation*}
 \mathcal B
 \equiv\frac{\alpha}{4\pi^2}
 \int\dd\Omega\,
 \left|
 \sum_i\sigma_i^{\rm io}Q_i
 \frac{\boldsymbol\beta_i^\perp}
 {1-\mathbf n\cdot\boldsymbol\beta_i}
 \right|^2\ge0.
\end{equation*}
于是未分辨真实光子在
$\lambda_{\rm IR}<\hbar\omega<\Delta E_{\rm det}$ 内的概率权重为
\begin{equation*}
 \mathcal R(\lambda_{\rm IR},\Delta E_{\rm det})
 =\mathcal B\int_{\lambda_{\rm IR}/\hbar}^{\Delta E_{\rm det}/\hbar}
 \frac{\dd\omega}{\omega}
 =\mathcal B\ln\frac{\Delta E_{\rm det}}{\lambda_{\rm IR}}.
\end{equation*}
这里 $\lambda_{\rm IR}$ 是计算用调节量，而 $\Delta E_{\rm det}$ 是实验把“有光子”和“无可分辨光子”划开的真实分辨率，两者物理地位不同。

为了把虚圈中的同一对数看清，引入任意软--硬分隔能量 $E_s$，满足所用软展开在
$\hbar\omega<E_s$ 内成立。由前一推导中的割线关系，虚软圈的红外部分具有同一 $\mathcal B\,\dd\omega/\omega$ 核而符号相反，因此可写成
\begin{equation*}
 2\Re\frac{\delta_{\rm virt}^{\rm soft}\mathcal M_0}{\mathcal M_0}
 =-\mathcal B\ln\frac{E_s}{\lambda_{\rm IR}}
 +C_{\rm virt}(E_s),
\end{equation*}
其中 $C_{\rm virt}(E_s)$ 在 $\lambda_{\rm IR}\to0$ 时有限。把对数拆成
\begin{align*}
 -\mathcal B\ln\frac{E_s}{\lambda_{\rm IR}}
 &=-\mathcal B\ln\frac{\Delta E_{\rm det}}{\lambda_{\rm IR}}
   -\mathcal B\ln\frac{E_s}{\Delta E_{\rm det}},
\end{align*}
就得到正文\eqnrefs{eq:virtual-real-ir-relation-dp10}，其中
\[
 C_{\rm IR\,finite}
 =C_{\rm virt}(E_s)-\mathcal B\ln\frac{E_s}{\Delta E_{\rm det}}.
\]

现在展开到 $O(\alpha)$。无真实可分辨光子的虚修正概率为
\begin{align*}
 |\mathcal M_0+\delta_{\rm virt}\mathcal M_0|^2
 &=|\mathcal M_0|^2
 \left[1+2\Re\frac{\delta_{\rm virt}\mathcal M_0}{\mathcal M_0}\right]
 +O(\alpha^2),
\end{align*}
而一个未分辨真实软光子的概率为
$|\mathcal M_0|^2\mathcal R+O(\alpha^2)$。两类末态在 Fock 空间正交，所以包容概率相加而不是振幅相加：
\begin{align*}
 P_{\rm incl}
 &=|\mathcal M_0|^2
 \left[
 1-\mathcal R+C_{\rm IR\,finite}+\mathcal R
 \right]+O(\alpha^2)\\
 &=|\mathcal M_0|^2
 \left[1+C_{\rm IR\,finite}\right]+O(\alpha^2).
\end{align*}
因此所有 $\ln\lambda_{\rm IR}$ 严格抵消，取
$\lambda_{\rm IR}\to0$ 得到正文\eqnrefs{eq:bloch-nordsieck-cancellation-dp10}。留下的
$\Delta E_{\rm det}$ 依赖并不应消失：改变探测能量分辨率就改变了“哪些软光子被计入同一个包容事件”的测量定义。
\end{derivation}

\begin{derivation}[软共线分母中的质量与角度尺度]
取电子四动量和光子四动量

\begin{equation*}
 p^\mu=(E/c,\mathbf p),
 \qquad
 k^\mu=(\hbar\omega/c,\hbar\omega\mathbf n/c),
\end{equation*}
其中 $|\mathbf n|=1$，$\mathbf p\cdot\mathbf n=p\cos\theta$。于是
\begin{equation*}
 p\!\cdot k
 =\frac{E\hbar\omega}{c^2}
 -\frac{p\hbar\omega}{c}\cos\theta
 =\frac{E\hbar\omega}{c^2}(1-\beta_{\rm rel}\cos\theta),
\end{equation*}
其中 $\beta_{\rm rel}=pc/E$。小角时
\begin{equation*}
 \cos\theta=1-\frac{\theta^2}{2}+\order(\theta^4).
\end{equation*}
质量壳关系给
\begin{equation*}
 \beta_{\rm rel}^2=1-\frac{m_{\mathrm e}^2c^4}{E^2},
\end{equation*}
令 $x\equiv m_{\mathrm e}^2c^4/E^2$。质量壳关系给 $\beta_{\rm rel}=\sqrt{1-x}$，所以
\begin{align*}
1-\beta_{\rm rel}
&=\frac{x}{1+\sqrt{1-x}}\\
&=\frac{x}{2}+O(x^2)
=\frac{m_{\mathrm e}^2c^4}{2E^2}
+O\!\left(\frac{m_{\mathrm e}^4c^8}{E^4}\right).
\end{align*}
保留 $\theta^2$ 和 $m_{\mathrm e}^2c^4/E^2$ 的最低非零阶，得到\eqnrefs{eq:soft-collinear-denominator-dp10}。这也显示电子质量如何把严格共线分母从零抬升为有限值。
\end{derivation}

\begin{derivation}[可区分费米子的 Yukawa 交换振幅]
从相互作用

\begin{equation*}
\mathcal L_{\rm int}
=-\sqrt{\hbar c}\,\phi
\left(y_A\bar\psi_A\psi_A+y_B\bar\psi_B\psi_B\right)
\end{equation*}
出发。二阶 Dyson 项含两个 Yukawa 顶角。每个顶角内的费米子双线性型都保持种类，因此外线收缩必须把入射 $A$ 与出射 $A$ 相连、入射 $B$ 与出射 $B$ 相连。于是两顶角之间的标量收缩携带
\begin{equation*}
q=p_1-p_3=p_4-p_2,
\qquad q^2=t.
\end{equation*}
若交换 $p_3\leftrightarrow p_4$，则末态种类也被交换，得到的是不同的外态标签而非同一物理末态的第二种 Wick 收缩，所以可区分粒子没有 $u$-道干涉。

逐个写约化因子，
\begin{align*}
\ii\mathcal M_t
&=\bigl[\bar u_A(p_3)(-\ii y_A)u_A(p_1)\bigr]
\frac{\ii}{t-m_Y^2c^2+\ii0^+}
\bigl[\bar u_B(p_4)(-\ii y_B)u_B(p_2)\bigr],\\
\mathcal M_t
&=\frac{y_Ay_B}{t-m_Y^2c^2}
[\bar u_A(p_3)u_A(p_1)]
[\bar u_B(p_4)u_B(p_2)].
\end{align*}
两个双线性型各有一阶动量量纲，正好抵消传播子的二阶动量量纲，因此约化振幅无量纲，与全书散射归一化一致。
\end{derivation}

\begin{derivation}[相同费米子交换通道的干涉]
干涉项来自 $t$、$u$ 两种振幅的交叉乘积。对初态两自旋平均并对末态求和后，利用完备关系

\begin{equation*}
\sum_su(p,s)\bar u(p,s)=\slashed p+mc,
\end{equation*}
可把四个双线性型首尾连接成
\begin{equation*}
\mathcal I_{tu}
=\frac14\Tr\!\left[
(\slashed p_3+mc)(\slashed p_1+mc)
(\slashed p_4+mc)(\slashed p_2+mc)
\right].
\end{equation*}
先按斜线记号的个数展开。含一个或三个 gamma 的迹为零，只剩零、二、四个斜线记号：
\begin{align*}
4\mathcal I_{tu}
={}&\Tr(\slashed p_3\slashed p_1\slashed p_4\slashed p_2)\\
&+m^2c^2\Bigl[
\Tr(\slashed p_3\slashed p_1)
+\Tr(\slashed p_3\slashed p_4)
+\Tr(\slashed p_3\slashed p_2)\\
&\hspace{3.2cm}
+\Tr(\slashed p_1\slashed p_4)
+\Tr(\slashed p_1\slashed p_2)
+\Tr(\slashed p_4\slashed p_2)
\Bigr]
+4m^4c^4.
\end{align*}
使用
\begin{align*}
\Tr(\slashed a\slashed b)&=4a\!\cdot b,\\
\Tr(\slashed a\slashed b\slashed c\slashed d)
&=4\bigl[(a\!\cdot b)(c\!\cdot d)
-(a\!\cdot c)(b\!\cdot d)
+(a\!\cdot d)(b\!\cdot c)\bigr],
\end{align*}
得到所需点积形式。下一步不能直接写“用 Mandelstam 变量整理”，而逐个代入等质量关系
\begin{align*}
p_1\!\cdot p_2&=p_3\!\cdot p_4=\frac{s-2m^2c^2}{2},\\
p_1\!\cdot p_3&=p_2\!\cdot p_4=m^2c^2-\frac t2,\\
p_1\!\cdot p_4&=p_2\!\cdot p_3=m^2c^2-\frac u2,
\end{align*}
并用
\begin{equation*}
s+t+u=4m^2c^2
\quad\Longrightarrow\quad
u=4m^2c^2-s-t.
\end{equation*}
四-gamma 部分先化为
\begin{align*}
&(p_3\!\cdot p_1)(p_4\!\cdot p_2)
-(p_3\!\cdot p_4)(p_1\!\cdot p_2)
+(p_3\!\cdot p_2)(p_1\!\cdot p_4)\\
&=\left(m^2c^2-\frac t2\right)^2
-\left(\frac{s-2m^2c^2}{2}\right)^2
+\left(m^2c^2-\frac u2\right)^2.
\end{align*}
六个二-gamma 项的点积和为
\begin{equation*}
\sum_{i<j}p_i\!\cdot p_j
=s+2m^2c^2,
\end{equation*}
其中最后一步由上述三个 Mandelstam 配对逐项相加得到。把 $u=4m^2c^2-s-t$ 代回并展开平方，所有 $u$ 依赖消去，最终得到
\begin{equation*}
\mathcal I_{tu}
=\frac12\left(4m^2c^2s-4m^2c^2t+st+t^2\right),
\end{equation*}
即目标式。质量维数检查给 $[\mathcal I_{tu}]=({\rm momentum})^4$，与两个传播子分母的总四阶量纲相消。
\end{derivation}

\begin{derivation}[Yukawa 散射的自旋平均]
取 $m_A=m_B=m$、$y_A=y_B=y$。由\eqnrefs{eq:yukawa-distinguishable-M-si-r10}，对两个初态自旋各平均 $1/2$，并利用完备关系

\begin{equation*}
 \sum_su(p,s)\bar u(p,s)=\slashed p+mc,
\end{equation*}
得到
\begin{equation*}
 \overline{|\mathcal M_t|^2}
 =\frac{y^4}{(t-m_Y^2c^2)^2}\frac14
 \Tr[(\slashed p_3+mc)(\slashed p_1+mc)]
 \Tr[(\slashed p_4+mc)(\slashed p_2+mc)].
\end{equation*}
第一条迹逐项展开为
\begin{align*}
 \Tr[(\slashed p_3+mc)(\slashed p_1+mc)]
 &=\Tr(\slashed p_3\slashed p_1)+m^2c^2\Tr\id(4)\\
 &=4(p_3\!\cdot p_1+m^2c^2),
\end{align*}
其中单个 $\gamma$ 矩阵的迹为零。等质量质量壳条件给
\begin{equation*}
 t=(p_1-p_3)^2=2m^2c^2-2p_1\!\cdot p_3,
\end{equation*}
因此
\begin{equation*}
 p_1\!\cdot p_3=m^2c^2-\frac t2,
\end{equation*}
从而
\begin{equation*}
 \Tr[(\slashed p_3+mc)(\slashed p_1+mc)]
 =2(4m^2c^2-t).
\end{equation*}
第二条费米线对应
\begin{align*}
 \Tr[(\slashed p_4+mc)(\slashed p_2+mc)]
 &=4(p_4\!\cdot p_2+m^2c^2).
\end{align*}
由同一个动量转移
\[
 t=(p_2-p_4)^2=2m^2c^2-2p_2\!\cdot p_4
\]
得到
\[
 p_2\!\cdot p_4=m^2c^2-\frac t2,
\]
因此第二条迹也等于 $2(4m^2c^2-t)$。两条迹相乘并乘初态自旋平均因子 $1/4$ 后，
\begin{align*}
 \overline{|\mathcal M_t|^2}
 &=\frac14\,
 \frac{y^4}{(t-m_Y^2c^2)^2}
 [2(4m^2c^2-t)]^2,\\
 &=y^4\frac{(4m^2c^2-t)^2}{(t-m_Y^2c^2)^2}.
\end{align*}
$t$ 与 $m_Y^2c^2$ 都是动量平方，所以结果无量纲；这与约化散射振幅的归一化一致。
\end{derivation}

\begin{derivation}[Yukawa 散射的质心系角分布]

正文\eqnrefs{eq:yukawa-distinguishable-M2-si-r10} 已经给出 Lorentz 不变的自旋平均振幅平方。转到质心系后，用 $t=-2p_*^2(1-\cos\theta)$ 和二体相空间通式消去 Mandelstam 变量，即可把不变结果改写成直接可测的角分布~\eqnrefs{eq:yukawa-distinguishable-angular-si-r10}；这一阶段不再引入新的动力学近似。

质心系取
\begin{equation*}
 p_1^\mu=(E/c,\mathbf p),\qquad
 p_2^\mu=(E/c,-\mathbf p),
\end{equation*}
且弹性散射满足 $|\mathbf p_3|=|\mathbf p_1|=p$。于是
\begin{equation*}
 s=(p_1+p_2)^2=\frac{4E^2}{c^2}=4(m^2c^2+p^2).
\end{equation*}
令 $\theta$ 为 $\mathbf p_1$ 与 $\mathbf p_3$ 的夹角，则
\begin{align*}
 t&=(p_1-p_3)^2\\
  &=-|\mathbf p_1-\mathbf p_3|^2\\
  &=-2p^2(1-\cos\theta).
\end{align*}
可区分二体弹性散射的 SI 主式为
\begin{equation*}
 \frac{\dd\sigma}{\dd\Omega}
 =\frac{\hbar^2}{64\pi^2s}
 \frac{|\mathbf p_f|}{|\mathbf p_i|}
 \overline{|\mathcal M|^2}.
\end{equation*}
这里动量比为一。把\eqnrefs{eq:yukawa-distinguishable-M2-si-r10} 与上述 $t$ 代入，即得目标角分布。截面量纲来自 $\hbar^2/s$，因为 $[s]=({\rm kg\,m\,s^{-1}})^2$，故 $[\hbar^2/s]={\rm m^2}$。
\end{derivation}

\begin{derivation}[费米子交换为何产生相对负号]
费米场的 Wick 收缩在交换两条外部费米算符时产生一次奇置换。用外态定义

\begin{equation*}
 |p_1s_1,p_2s_2\rangle
 =b_{s_1}^\dagger(p_1)b_{s_2}^\dagger(p_2)|0\rangle
\end{equation*}
和
\begin{equation*}
 \{b_r(\mathbf p),b_s^\dagger(\mathbf p')\}
 =(2\pi\hbar)^3\delta_{rs}\delta^{(3)}(\mathbf p-\mathbf p')
\end{equation*}
交换末态两个产生算符需要一次反对易，因而第二种收缩相对第一种带负号。分别代入两个标量传播子
\begin{equation*}
 \frac{\ii}{t-m_Y^2c^2+\ii0^+},
 \qquad
 \frac{\ii}{u-m_Y^2c^2+\ii0^+},
\end{equation*}
并使用相同的两个 Yukawa 顶角 $-\ii y$，便得到\eqnrefs{eq:yukawa-identical-M-si-r10}。
\end{derivation}

\begin{derivation}[Yukawa 交换的低速势能]
协变归一化的正能 Dirac 旋量可写成

\begin{equation*}
 u(p)=\sqrt{\frac{E+mc^2}{c}}
 \begin{pmatrix}
 \xi\\[1mm]
 \dfrac{c\boldsymbol\sigma\cdot\mathbf p}{E+mc^2}\xi
 \end{pmatrix}.
\end{equation*}
当 $|\mathbf p|/(mc)\ll1$ 时，
\begin{equation*}
 E=mc^2+\frac{p^2}{2m}+\order\!\left(\frac{p^4}{m^3c^2}\right),
\end{equation*}
因此大分量给
\begin{equation*}
 \bar u(p')u(p)
 =2mc\,\xi'^\dagger\xi
 +\order\!\left(\frac{p^2}{mc}\right).
\end{equation*}
Mandelstam 变量始终满足严格恒等式
\begin{equation*}
 t=-\mathbf q^2+(q^0)^2.
\end{equation*}
定义静态非相对论振幅为同时取 $|\mathbf p|/(mc)\to0$ 与 $q^0/|\mathbf q|\to0$ 的极限。利用上面的旋量展开，便有
\begin{align*}
 \mathcal M_t^{\rm NR}(\mathbf q)
 &\equiv
 \lim_{\substack{|\mathbf p|/(mc)\to0\\ q^0/|\mathbf q|\to0}}
 \mathcal M_t\\
 &=-\frac{4m^2c^2y^2}{\mathbf q^2+m_Y^2c^2}.
\end{align*}
有限速度与有限能量转移的修正由前式中 $O(p^2/m^2c^2)$ 的旋量余项以及精确分母中的 $(q^0)^2$ 项逐项恢复。

Schr\"odinger 第一 Born 近似与上述约化振幅的归一化关系给出
\begin{equation*}
 \widetilde V_Y(\mathbf q)
 =-\frac{4\pi\alpha_Y\hbar^3c}
 {\mathbf q^2+m_Y^2c^2},
 \qquad \alpha_Y=\frac{y^2}{4\pi}.
\end{equation*}
再使用后续推导中得到的三维 Fourier 积分
\begin{equation*}
 \int\frac{\dd^3 q}{(2\pi\hbar)^3}
 \frac{\ee^{\ii\mathbf q\cdot\mathbf r/\hbar}}
 {\mathbf q^2+m_Y^2c^2}
 =\frac{\ee^{-m_Ycr/\hbar}}{4\pi\hbar^2r},
\end{equation*}
得到
\begin{equation*}
 V_Y(r)=-\alpha_Y\hbar c\frac{\ee^{-m_Ycr/\hbar}}{r}.
\end{equation*}
指数中的无量纲组合是 $m_Ycr/\hbar=r/\lambda_Y$，因此 $\lambda_Y=\hbar/(m_Yc)$ 是媒介粒子的 Compton 作用程。$m_Y\to0$ 恢复 $1/r$ 长程极限，$r\gg\lambda_Y$ 时相互作用指数抑制。
\end{derivation}

\begin{derivation}[低速电磁电流与 Coulomb 势]
低速势能由电磁交换振幅决定，因而需要先计算电流的低速展开。使用前文归一化的正能 Dirac 旋量，对质量为 $m$ 的外线记
\begin{equation*}
 D_{\mathbf p}\equiv \frac{E(\mathbf p)}{c}+mc,
 \qquad
 N_{\mathbf p}\equiv\sqrt{D_{\mathbf p}},
\end{equation*}
则正文\eqnrefs{eq:explicit-u} 可写成
\begin{equation*}
 u_s(p)=N_{\mathbf p}
 \begin{pmatrix}
 \xi_s\\[1mm]
 \dfrac{\boldsymbol\sigma_{\rm P}\cdot\mathbf p}{D_{\mathbf p}}\xi_s
 \end{pmatrix},
 \qquad
 \xi_{s'}^\dagger\xi_s=\delta_{s's}.
\end{equation*}
由精确色散
$E(\mathbf p)=\sqrt{m^2c^4+c^2\mathbf p^2}$，在
$|\mathbf p|/(mc)\ll1$ 时
\begin{align*}
 \frac{E(\mathbf p)}{c}
 &=mc+\frac{\mathbf p^2}{2mc}
 +O\!\left(\frac{|\mathbf p|^4}{m^3c^3}\right),\\
 D_{\mathbf p}
 &=2mc+\frac{\mathbf p^2}{2mc}
 +O\!\left(\frac{|\mathbf p|^4}{m^3c^3}\right),\\
 N_{\mathbf p}
 &=\sqrt{2mc}\left[
 1+\frac{\mathbf p^2}{8m^2c^2}
 +O\!\left(\frac{|\mathbf p|^4}{m^4c^4}\right)
 \right].
\end{align*}

时间分量的 Dirac 电流为
\begin{align*}
 \bar u_{s'}(p')\gamma^0u_s(p)
 &=u_{s'}^\dagger(p')u_s(p)\\
 &=N_{\mathbf p'}N_{\mathbf p}\,
 \xi_{s'}^\dagger\left[
 I_2+
 \frac{(\boldsymbol\sigma_{\rm P}\cdot\mathbf p')
       (\boldsymbol\sigma_{\rm P}\cdot\mathbf p)}
 {D_{\mathbf p'}D_{\mathbf p}}
 \right]\xi_s.
\end{align*}
第二项已经含两个小动量；同时
$N_{\mathbf p'}N_{\mathbf p}=2mc+O(|\mathbf p|^2/(mc),|\mathbf p'|^2/(mc))$，故
\begin{equation*}
 \boxed{
 \bar u_{s'}(p')\gamma^0u_s(p)
 =2mc\,\delta_{s's}
 +O\!\left(\frac{|\mathbf p|^2+|\mathbf p'|^2}{mc}\right).}
\end{equation*}

空间分量不能仅靠量纲判断。由
$\bar u\gamma^iu=u^\dagger\alpha^iu$ 与
$\alpha^i=\begin{psmallmatrix}0&\sigma_{\rm P}^i\\
\sigma_{\rm P}^i&0\end{psmallmatrix}$，有
\begin{align*}
 \bar u_{s'}(p')\gamma^iu_s(p)
 ={}&N_{\mathbf p'}N_{\mathbf p}\,
 \xi_{s'}^\dagger\left[
 \frac{\sigma_{\rm P}^i(\boldsymbol\sigma_{\rm P}\cdot\mathbf p)}{D_{\mathbf p}}
 +\frac{(\boldsymbol\sigma_{\rm P}\cdot\mathbf p')\sigma_{\rm P}^i}{D_{\mathbf p'}}
 \right]\xi_s.
\end{align*}
利用
$\sigma_{\rm P}^i\sigma_{\rm P}^j
=\delta^{ij}I_2+\ii\epsilon^{ijk}\sigma_{\rm P}^k$，领先阶成为
\begin{equation*}
 \boxed{
 \bar u_{s'}(p')\gamma^iu_s(p)
 =\xi_{s'}^\dagger\left[
 (p+p')^i I_2
 +\ii\epsilon^{ijk}(p-p')^j\sigma_{\rm P}^k
 \right]\xi_s
 +O\!\left(\frac{|\mathbf p|^3+|\mathbf p'|^3}{m^2c^2}\right).}
\end{equation*}
所以时间分量为 $O(mc)$，空间分量为 $O(|\mathbf p|)$；这才是静态极限中电荷密度压倒三维电流的数学原因。

对两条外费米子线分别应用这一结果，并取不翻转自旋通道，正文\eqnrefs{eq:qed-reduced-exchange-r13} 给出
\begin{align*}
 \mathcal M_\gamma
 &=\frac{g_1g_2}{q^2}
 J_1^\mu J_{2\mu}\\
 &=\frac{g_1g_2}{q^2}
 \left[J_1^0J_2^0-\mathbf J_1\cdot\mathbf J_2\right]\\
 &=\frac{4m_1m_2c^2g_1g_2}{q^2}
 +O\!\left(\frac{|\mathbf p|^2}{m_1m_2c^2}\right).
\end{align*}
近静态交换再给
\begin{equation*}
 q^2=(q^0)^2-\mathbf q^2
 =-\mathbf q^2\left[
 1-\frac{(q^0)^2}{\mathbf q^2}
 \right],
\end{equation*}
故
\begin{equation*}
 \mathcal M_\gamma^{\rm NR}(\mathbf q)
 =-\frac{4m_1m_2c^2g_1g_2}{\mathbf q^2}
 +O\!\left(
 \frac{|\mathbf p|^2}{\mathbf q^2},
 \frac{m_1m_2c^2g_1g_2(q^0)^2}{|\mathbf q|^4}
 \right).
\end{equation*}
把领先项代入前文已经从相对论与 Schrödinger 截面归一化推得的匹配式
\eqnrefs{eq:relativistic-to-nr-potential-matching-dp75}：
\begin{align*}
 \widetilde V_C(\mathbf q)
 &=-\frac{\hbar^3}{4m_1m_2c}
 \mathcal M_\gamma^{\rm NR}(\mathbf q)\\
 &=\frac{\hbar^3c}{\mathbf q^2}
 \frac{q_1q_2}{\varepsilon_0\hbar c}\\
 &=\boxed{\frac{q_1q_2\hbar^2}
 {\varepsilon_0\mathbf q^2}},
\end{align*}
即正文\eqnrefs{eq:coulomb-Vq-r13}。随后只剩三维 Fourier 变换
\begin{equation*}
 \int\frac{\dd^3q}{(2\pi\hbar)^3}
 \frac{\ee^{\ii\mathbf q\cdot\mathbf r/\hbar}}{\mathbf q^2}
 =\frac{1}{4\pi\hbar^2r},
\end{equation*}
因此
$V_C(r)=q_1q_2/(4\pi\varepsilon_0r)$。这最后一步是标准 Fourier 积分；真正的物理内容在前面的三步：Dirac 电流低速展开、近静态传播子极限以及相对论振幅到势核的归一化匹配。
\end{derivation}

\begin{derivation}[电子与缪子的散射振幅平方]
取外态

\begin{equation*}
 e^-(p_1,s_1)+\mu^-(p_2,s_2)
 \to e^-(p_3,s_3)+\mu^-(p_4,s_4),
\end{equation*}
则 $q=p_1-p_3$、$q^2=t$，振幅为
\begin{align*}
 \ii\mathcal M_t
 &=(-\ii g_{\rm em})\bar u_e(p_3)\gamma^\mu u_e(p_1)
 \frac{-\ii\eta_{\mu\nu}}{t}
 (-\ii g_{\rm em})\bar u_\mu(p_4)\gamma^\nu u_\mu(p_2),\\
 \mathcal M_t
 &=\frac{g_{\rm em}^2}{t}
 [\bar u_3\gamma^\mu u_1]
 [\bar u_4\gamma_\mu u_2].
 \end{align*}
未偏振实验对电子和 $\mu$ 子初态各平均两种自旋，所以
\begin{align*}
 \overline{|\mathcal M_t|^2}
 =&\frac{g_{\rm em}^4}{t^2}\frac14
 \Tr[(\slashed p_3+m_{\mathrm e} c)\gamma^\mu(\slashed p_1+m_{\mathrm e} c)\gamma^\nu]\nonumber\\
 &\hspace{2.4cm}\times
 \Tr[(\slashed p_4+m_\mu c)\gamma_\mu(\slashed p_2+m_\mu c)\gamma_\nu].
 \end{align*}
使用 \eqnrefs{eq:qed-basic-tensor-si-r9}，先把两条费米子线的迹分别写成秩二张量：
\begin{align*}
 L_\ee^{\mu\nu}
 &=4\left[
 p_3^\mu p_1^\nu+p_3^\nu p_1^\mu
 -(p_3\!\cdot p_1-m_{\mathrm e}^2c^2)\eta^{\mu\nu}
 \right],\\
 L_\mu^{\mu\nu}
 &=4\left[
 p_4^\mu p_2^\nu+p_4^\nu p_2^\mu
 -(p_4\!\cdot p_2-m_\mu^2c^2)\eta^{\mu\nu}
 \right].
 \end{align*}
于是 \eqnrefs{eq:emu-M2-trace} 变成
\begin{align*}
 \overline{|\mathcal M_t|^2}
 &=\frac{g_{\rm em}^4}{4t^2}L_\ee^{\mu\nu}L_{\mu,\mu\nu},\\
 &=\frac{g_{\rm em}^4}{t^2}\left\{
 8\left[(p_1\!\cdot p_2)(p_3\!\cdot p_4)
 +(p_1\!\cdot p_4)(p_2\!\cdot p_3)\right]
 +4t(m_{\mathrm e}^2+m_\mu^2)c^2
 \right\}.
 \end{align*}
这里第二行只做了 Lorentz 指标收缩。再由
\begin{align*}
 2p_1\!\cdot p_2=2p_3\!\cdot p_4
 &=s-(m_{\mathrm e}^2+m_\mu^2)c^2\equiv S_0,\\
 -2p_1\!\cdot p_4=-2p_2\!\cdot p_3
 &=u-(m_{\mathrm e}^2+m_\mu^2)c^2\equiv U_0,
\end{align*}
得到
\begin{align*}
 \overline{|\mathcal M_t|^2}
 =\frac{2g_{\rm em}^4}{t^2}
 \left[
 S_0^2+U_0^2+2t(m_{\mathrm e}^2+m_\mu^2)c^2
 \right],
\end{align*}
即正文\eqnrefs{eq:emu-M2-massive-invariants-dp49}。


\end{derivation}

\begin{derivation}[Møller 散射的交换干涉与角分布]
取入射动量为 $p_1,p_2$，出射动量为 $p_3,p_4$。由于两条末态电子不可区分，有两种 Wick 收缩：

\begin{align*}
 \mathcal M_t
 &=\frac{g_{\rm em}^2}{t}
 [\bar u_3\gamma^\mu u_1]
 [\bar u_4\gamma_\mu u_2],\\
 \mathcal M_u
 &=\frac{g_{\rm em}^2}{u}
 [\bar u_4\gamma^\mu u_1]
 [\bar u_3\gamma_\mu u_2].
\end{align*}
交换两个外部费米子算符需要一次奇置换，因此
\begin{equation*}
 {
 \mathcal M_{\mathrm{Moller}}=\mathcal M_t-\mathcal M_u.}
 \end{equation*}
平方以后
\begin{equation*}
 |\mathcal M_t-\mathcal M_u|^2
 =|\mathcal M_t|^2+|\mathcal M_u|^2
 -2\Re(\mathcal M_t\mathcal M_u^*).
\end{equation*}
在 $s\gg m_{\mathrm e}^2c^2$ 极限，两个平方项就是前面单一交换迹的两种外部标签分配：
\begin{align*}
 \overline{|\mathcal M_t|^2}
 &=2g_{\rm em}^4\frac{s^2+u^2}{t^2},\\
 \overline{|\mathcal M_u|^2}
 &=2g_{\rm em}^4\frac{s^2+t^2}{u^2}.
\end{align*}
干涉项需要单独计算。把四条自旋求和首尾闭合成一个迹，得到
\begin{align*}
 \overline{-2\Re(\mathcal M_t\mathcal M_u^*)}
 &=-\frac{g_{\rm em}^4}{2tu}
 \Tr\!\left[
 \slashed p_3\gamma^\mu\slashed p_1\gamma^\nu
 \slashed p_4\gamma_\mu\slashed p_2\gamma_\nu
 \right].
 \end{align*}
使用前面由 Clifford 代数建立的
\eqnrefs{eq:gamma-two-sandwich-r10,eq:gamma-three-sandwich-r10}，迹依次化为
\begin{align*}
 \mathcal I_{tu}
 &\equiv
 \Tr[\slashed p_3\gamma^\mu\slashed p_1\gamma^\nu
 \slashed p_4\gamma_\mu\slashed p_2\gamma_\nu],\\
 &=-2\Tr[\slashed p_3\gamma^\mu\slashed p_1\slashed p_2
 \gamma_\mu\slashed p_4],\\
 &=-32(p_1\!\cdot p_2)(p_3\!\cdot p_4)
 =-8s^2.
\end{align*}
因此
\begin{equation*}
 \overline{-2\Re(\mathcal M_t\mathcal M_u^*)}
 =4g_{\rm em}^4\frac{s^2}{tu}.
 \end{equation*}
因此
\begin{equation*}
 {
 \overline{|\mathcal M_{\mathrm{Moller}}|^2}
 =2g_{\rm em}^4\left[
 \frac{s^2+u^2}{t^2}
 +\frac{s^2+t^2}{u^2}
 +\frac{2s^2}{tu}
 \right].}
 \end{equation*}
若 $\theta$ 在全 $0\le\theta\le\pi$ 上标记“第 3 条电子”的方向，则 $(p_3,p_4)$ 与 $(p_4,p_3)$ 被重复计数，需要 $1/2!$。使用
\begin{equation*}
 t=-\frac s2(1-\cos\theta),
 \qquad
 u=-\frac s2(1+\cos\theta)
\end{equation*}
可化为
\begin{equation*}
 {
 \frac{\dd\sigma_{\mathrm{Moller}}}{\dd\Omega}
 =\frac{\hbar^2\alpha^2}{2s}
 \frac{(3+\cos^2\theta)^2}{\sin^4\theta}.}
 \end{equation*}
若只在一个不重复的半球区域定义末态，则不再额外除 $2!$，两种约定必须二选一。
\end{derivation}

\begin{derivation}[Bhabha 散射的通道干涉与角分布]
Bhabha 散射同时有湮灭的 $s$ 通道与散射的 $t$ 通道：

\begin{align*}
 \mathcal M_s
 &=\frac{g_{\rm em}^2}{s}
 [\bar v_2\gamma^\mu u_1]
 [\bar u_3\gamma_\mu v_4],\\
 \mathcal M_t
 &=-\frac{g_{\rm em}^2}{t}
 [\bar u_3\gamma^\mu u_1]
 [\bar v_2\gamma_\mu v_4].
 \end{align*}
相对号由把同一个 Wick 收缩次序写成标准外部旋量次序时确定。对无偏振、$s\gg m_{\mathrm e}^2c^2$，两个平方项分别复用湮灭与交换的迹：
\begin{align*}
 \overline{|\mathcal M_s|^2}
 &=2g_{\rm em}^4\frac{u^2+t^2}{s^2},\\
 \overline{|\mathcal M_t|^2}
 &=2g_{\rm em}^4\frac{u^2+s^2}{t^2}.
\end{align*}
交叉项把四条外部自旋求和s 连成
\begin{align*}
 2\Re\,\overline{\mathcal M_s\mathcal M_t^*}
 &=-\frac{g_{\rm em}^4}{2st}
 \Tr\!\left[
 \slashed p_2\gamma^\mu\slashed p_1\gamma^\nu
 \slashed p_3\gamma_\mu\slashed p_4\gamma_\nu
 \right].
 \end{align*}
使用 \eqnrefs{eq:gamma-two-sandwich-r10,eq:gamma-three-sandwich-r10} 的 $\gamma$ 矩阵夹心恒等式，
\begin{align*}
 \mathcal I_{st}
 &=-32(p_1\!\cdot p_4)(p_2\!\cdot p_3),\\
 &=-8u^2,
\end{align*}
所以
\begin{equation*}
 2\Re\,\overline{\mathcal M_s\mathcal M_t^*}
 =4g_{\rm em}^4\frac{u^2}{st}.
\end{equation*}
三项相加为
\begin{equation*}
 {
 \overline{|\mathcal M_{\rm Bhabha}|^2}
 =2g_{\rm em}^4\left[
 \frac{u^2+t^2}{s^2}
 +\frac{u^2+s^2}{t^2}
 +\frac{2u^2}{st}
 \right].}
 \end{equation*}
代入质心系的 $t,u$ 后
\begin{equation*}
 {
 \frac{\dd\sigma_{\rm Bhabha}}{\dd\Omega}
 =\frac{\hbar^2\alpha^2}{4s}
 \frac{(3+\cos^2\theta)^2}{(1-\cos\theta)^2}.}
 \end{equation*}
\end{derivation}

\begin{derivation}[散射角与动量转移的变量变换]
在质心系写

\begin{equation*}
 p_1^\mu=(E_1/c,\mathbf p_i),
 \qquad
 p_3^\mu=(E_3/c,\mathbf p_f),
\end{equation*}
并令 $\theta$ 为 $\mathbf p_i$ 与 $\mathbf p_f$ 的夹角。由 $t=(p_1-p_3)^2$，
\begin{align*}
 t
 &=p_1^2+p_3^2-2p_1\!\cdot p_3,\\
 &=(m_1^2+m_3^2)c^2
 -\frac{2E_1E_3}{c^2}
 +2|\mathbf p_i||\mathbf p_f|\cos\theta.
 \end{align*}
固定 $s$ 时，$E_1,E_3,|\mathbf p_i|,|\mathbf p_f|$ 都不随 $\theta$ 改变，因此
\begin{equation*}
 \frac{\dd t}{\dd\cos\theta}
 =2|\mathbf p_i||\mathbf p_f|.
\end{equation*}
轴对称两体散射有
\begin{equation*}
 \dd\Omega=2\pi\,\dd(\cos\theta),
 \qquad
 \frac{\dd\Omega}{\dd t}
 =\frac{\pi}{|\mathbf p_i||\mathbf p_f|}.
\end{equation*}
代入二体截面主式，
\begin{align*}
 \frac{\dd\sigma}{\dd t}
 &=\frac{\dd\sigma}{\dd\Omega}
 \frac{\dd\Omega}{\dd t},\\
 &=\frac{\hbar^2}{64\pi^2s}
 \frac{|\mathbf p_f|}{|\mathbf p_i|}
 \overline{|\mathcal M|^2}
 \frac{\pi}{|\mathbf p_i||\mathbf p_f|},\\
 &=\frac{\hbar^2\overline{|\mathcal M|^2}}
 {64\pi s|\mathbf p_i|^2}.
\end{align*}
再由 \eqnrefs{eq:cm-momenta-kallen-si-r9}
\begin{equation*}
 |\mathbf p_i|^2
 =\frac{\Lambda_K(s,m_1^2c^2,m_2^2c^2)}{4s},
\end{equation*}
得到
\begin{equation*}
 {
 \frac{\dd\sigma}{\dd t}
 =\frac{\hbar^2\overline{|\mathcal M|^2}}
 {16\pi\Lambda_K(s,m_1^2c^2,m_2^2c^2)}.}
 \end{equation*}
由于 $\Lambda_K$ 具有动量四次方量纲，右端的量纲是面积除以动量平方，与 $\dd\sigma/\dd t$ 一致。
\end{derivation}

\begin{derivation}[电子湮灭产生缪子对的总截面]

正文\eqnrefs{eq:ee-mumu-M} 的非极化平方振幅由初、末态 Dirac 张量的缩并给出；二体相空间把它转成角分布，角积分再给总截面。考虑 $s\gg m_{\mathrm e}^2c^2$ 而允许 $s$ 接近 $4m_\mu^2c^2$，因此忽略初态电子质量而保留末态 $m_\mu$。

末态与初态的 Dirac 张量分别为
\begin{align*}
 L_{\mu}^{\rho\sigma}
 &=\Tr[(\slashed p_3+m_\mu c)\gamma^\rho
 (\slashed p_4-m_\mu c)\gamma^\sigma]\\
 &=4\left[
 p_3^\rho p_4^\sigma+p_3^\sigma p_4^\rho
 -(p_3\!\cdot p_4+m_\mu^2c^2)\eta^{\rho\sigma}
 \right],\\
 L_\ee^{\rho\sigma}
 &=4\left(
 p_2^\rho p_1^\sigma+p_2^\sigma p_1^\rho
 -p_1\!\cdot p_2\,\eta^{\rho\sigma}
 \right).
\end{align*}
因此
\begin{align*}
 \overline{|\mathcal M|^2}
 &=\frac{g_{\rm em}^4}{4s^2}L_{e,\rho\sigma}L_\mu^{\rho\sigma}\\
 &=\frac{8g_{\rm em}^4}{s^2}
 \left[
 (p_2\!\cdot p_3)(p_1\!\cdot p_4)
 +(p_2\!\cdot p_4)(p_1\!\cdot p_3)
 +(p_1\!\cdot p_2)m_\mu^2c^2
 \right].
\end{align*}
在质心系中，正文\eqnrefs{eq:mumu-beta-si-r10} 给出
\begin{align*}
 p_1\!\cdot p_2&=\frac{s}{2},\\
 p_1\!\cdot p_3&=\frac{s}{4}
 (1-\beta_{\rm rel}^{(\mu)}\cos\theta),\\
 p_1\!\cdot p_4&=\frac{s}{4}
 (1+\beta_{\rm rel}^{(\mu)}\cos\theta),
\end{align*}
其余两项由 $1\leftrightarrow2$ 对称性得到。再用
$4m_\mu^2c^2/s=1-(\beta_{\rm rel}^{(\mu)})^2$，整理为正文\eqnrefs{eq:mumu-M2-massive-angle-si-r10}。

二体截面公式给
\begin{equation*}
 \frac{\dd\sigma}{\dd\Omega}
 =\frac{\hbar^2}{64\pi^2s}
 \frac{|\mathbf p_f|}{|\mathbf p_i|}
 \overline{|\mathcal M|^2}.
\end{equation*}
由 $|\mathbf p_f|/|\mathbf p_i|=\beta_{\rm rel}^{(\mu)}$ 与
$g_{\rm em}^2=4\pi\alpha$，得到正文\eqnrefs{eq:mumu-dsigma-massive-si-r10}。

最后令 $x=\cos\theta$，$\dd\Omega=2\pi\dd x$，则
\begin{align*}
 \sigma
 &=\frac{\hbar^2\alpha^2\beta_{\rm rel}^{(\mu)}}{4s}
 2\pi\int_{-1}^{1}\dd x\,
 \left[2-(\beta_{\rm rel}^{(\mu)})^2
 +(\beta_{\rm rel}^{(\mu)})^2x^2\right]\\
 &=\frac{\pi\hbar^2\alpha^2\beta_{\rm rel}^{(\mu)}}{s}
 \left(2-\frac{2}{3}(\beta_{\rm rel}^{(\mu)})^2\right).
\end{align*}
代入
$(\beta_{\rm rel}^{(\mu)})^2=1-4m_\mu^2c^2/s$，得到正文\eqnrefs{eq:mumu-total-massive-si-r10}。
\end{derivation}

\begin{derivation}[电子极点的在壳重整化条件]
正文把极点问题化成一个真正的一维标量问题。固定未来指向单位类时四矢量 $u^\mu$，令
\begin{equation*}
 p^\mu=\lambda u^\mu,
 \qquad
 P_+(u)=\frac{\id(4)+\slashed u}{2}.
\end{equation*}
因为 $u^2=1$，Clifford 关系给
\begin{equation*}
 (\slashed u)^2=\id(4),
 \qquad
 P_+\slashed u=\slashed uP_+=P_+.
\end{equation*}
因此自由逆传播子投影到正能子空间后只有一个标量因子：
\begin{align*}
 P_+(\slashed p-m_{\mathrm e} c)P_+
 &=P_+(\lambda\slashed u-m_{\mathrm e} c)P_+\\
 &=(\lambda-m_{\mathrm e} c)P_+.
\end{align*}

一般 Lorentz 协变、宇称偶的电子自能可写成
\begin{equation*}
 \Sigma_R(p)=A_R(p^2)\slashed p+B_R(p^2)m_{\mathrm e} c,
\end{equation*}
其中 $A_R,B_R$ 是标量函数。沿 $p=\lambda u$ 投影：
\begin{align*}
 P_+\Sigma_R(\lambda u)P_+
 &=\left[\lambda A_R(\lambda^2)+m_{\mathrm e} c\,B_R(\lambda^2)\right]P_+.
\end{align*}
这说明正文定义的
\begin{equation*}
 \sigma_R(\lambda)
 =\frac12\Tr\!\left[P_+\Sigma_R(\lambda u)\right]
\end{equation*}
确实就是正能二维子空间上的自能本征值，因为 $\Tr P_+=2$。于是完整逆传播子满足
\begin{equation*}
 P_+S_R^{-1}(\lambda u)P_+
 =d_+(\lambda)P_+,
\end{equation*}
其中
\begin{equation*}
 d_+(\lambda)
 \equiv\lambda-m_{\mathrm e} c-\sigma_R(\lambda).
\end{equation*}

物理单电子极点位于 $\lambda=m_{\mathrm e} c$ 的意思是正能投影逆传播子在那里有零本征值：
\begin{equation*}
 d_+(m_{\mathrm e} c)=0.
\end{equation*}
代入定义即得
\begin{equation*}
 \sigma_R(m_{\mathrm e} c)=0.
\end{equation*}
这一步只固定零点位置，也就是物理质量；还没有固定极点留数。

再把 $d_+$ 在零点附近作普通一元 Taylor 展开。令
\begin{equation*}
 \delta\lambda\equiv\lambda-m_{\mathrm e} c,
\end{equation*}
则
\begin{align*}
 d_+(m_{\mathrm e} c+\delta\lambda)
 &=d_+(m_{\mathrm e} c)
 +d_+'(m_{\mathrm e} c)\delta\lambda
 +\frac12d_+''(m_{\mathrm e} c)(\delta\lambda)^2+\cdots\\
 &=\left[1-\sigma_R'(m_{\mathrm e} c)\right]\delta\lambda
 +O((\delta\lambda)^2).
\end{align*}
因此正能极点部分的传播子为
\begin{equation*}
 S_R(\lambda u)P_+
 =\frac{\ii Z_{\rm pole}P_+}{\delta\lambda+\ii0^+}
 +S_{R,\mathrm{reg}}(\lambda;u)P_+,
\end{equation*}
其中正则项在 $\delta\lambda=0$ 处有限，而留数由
\begin{equation*}
 Z_{\rm pole}
 =\frac{1}{d_+'(m_{\mathrm e} c)}
 =\frac{1}{1-\sigma_R'(m_{\mathrm e} c)}
\end{equation*}
唯一确定。在壳方案把重整化电子场定义成单位极点留数，即
\begin{equation*}
 Z_{\rm pole}=1,
\end{equation*}
于是
\begin{equation*}
 \sigma_R'(m_{\mathrm e} c)=0.
\end{equation*}
与质量条件合起来正是正文\eqnrefs{eq:qed-onshell-selfenergy-conditions-dp53}。

这个推导同时说明了两条反项条件为什么独立：质量反项调节 $d_+$ 的零点位置，场强反项调节同一零点的一阶导数；它们分别控制 Laurent 展开的极点位置和留数。
\end{derivation}

\begin{derivation}[Compton 散射的自旋与偏振求和]
从\eqnrefs{eq:compton-M} 的两个电子交换通道出发，计算非极化振幅平方，再用电子静止系的二体运动学得到角分布。记入射电子、入射光子和出射光子的四动量分别为 $p,k,k'$，并定义
\begin{equation*}
 a\equiv p\cdot k,\qquad b\equiv p\cdot k',
\end{equation*}
并写
\begin{equation*}
 \mathcal K_{\rm BW}^{\mu\nu}
 =\gamma^\mu\frac{\slashed p+\slashed k+m_{\mathrm e} c}{2a}\gamma^\nu
 -\gamma^\nu\frac{\slashed p-\slashed k'+m_{\mathrm e} c}{2b}\gamma^\mu .
\end{equation*}
初态电子自旋和光子偏振各平均一次，末态求和。由于完整振幅已经满足 Ward 条件，偏振和中的规范项不贡献，因此
\begin{equation*}
 \overline{|\mathcal M_C|^2}
 =\frac{g_{\rm em}^4}{4}
 \Tr\!\left[(\slashed p'+m_{\mathrm e} c)\mathcal K_{\rm BW}^{\mu\nu}
 (\slashed p+m_{\mathrm e} c)\overline{\mathcal K}_{{\rm BW},\mu\nu}\right].
\end{equation*}
质量壳与动量守恒给出
\begin{equation*}
 k\cdot k'=a-b,\qquad
 p\cdot p'=m_{\mathrm e}^2c^2+a-b,
\end{equation*}
以及内部电子动量
\begin{equation*}
 r_s=p+k,\quad r_u=p-k',\quad
 r_s^2=m_{\mathrm e}^2c^2+2a,\quad r_u^2=m_{\mathrm e}^2c^2-2b.
\end{equation*}
把迹分成 $ss$、$uu$ 与两图干涉三部分。对 $ss$ 项，先用
\begin{equation*}
 \gamma^\nu(\slashed p+m_{\mathrm e} c)\gamma_\nu=-2\slashed p+4m_{\mathrm e} c
\end{equation*}
和
\begin{equation*}
 \slashed r\slashed p\slashed r=2(r\cdot p)\slashed r-r^2\slashed p
\end{equation*}
化简内部夹心，得到
\begin{align*}
 T_{ss}
 &=\Tr\!\left[(\slashed p'+m_{\mathrm e} c)\gamma^\mu(\slashed r_s+m_{\mathrm e} c)
 \gamma^\nu(\slashed p+m_{\mathrm e} c)\gamma_\nu(\slashed r_s+m_{\mathrm e} c)\gamma_\mu\right]\\
 &=32\left(m_{\mathrm e}^4c^4+m_{\mathrm e}^2c^2a+ab\right).
\end{align*}
连同平均因子和 $(2a)^{-2}$ 的传播子分母，
\begin{equation*}
 \overline{|\mathcal M_s|^2}
 =2g_{\rm em}^4\left(\frac ba+\frac{m_{\mathrm e}^2c^2}{a}+\frac{m_{\mathrm e}^4c^4}{a^2}\right).
\end{equation*}
$uu$ 项同样给出
\begin{equation*}
 \overline{|\mathcal M_u|^2}
 =2g_{\rm em}^4\left(\frac ab-\frac{m_{\mathrm e}^2c^2}{b}+\frac{m_{\mathrm e}^4c^4}{b^2}\right).
\end{equation*}
交叉迹在两种次序下相同，
\begin{equation*}
 T_{su}=T_{us}=16m_{\mathrm e}^2c^2(2m_{\mathrm e}^2c^2+a-b),
\end{equation*}
而正文振幅中的相对负号使干涉贡献成为
\begin{equation*}
 \overline{2\Re(\mathcal M_s\mathcal M_u^*)}
 =2g_{\rm em}^4\left[
 m_{\mathrm e}^2c^2\left(\frac1a-\frac1b\right)
 -\frac{2m_{\mathrm e}^4c^4}{ab}\right].
\end{equation*}
三部分相加并整理同次幂的 $1/a,1/b$，得到正文\eqnrefs{eq:compton-M2-unpolarized-dp73}。

电子静止系中
\begin{equation*}
 p^\mu=(m_{\mathrm e} c,\mathbf0),\qquad
 a=m_{\mathrm e}\hbar\omega,\qquad b=m_{\mathrm e}\hbar\omega'.
\end{equation*}
利用正文反冲关系~\eqnrefs{eq:compton-energy-ratio}，\eqnrefs{eq:compton-M2-unpolarized-dp73} 化为
\begin{equation*}
 \overline{|\mathcal M_C|^2}
 =2g_{\rm em}^4\left(
 \frac\omega{\omega'}+\frac{\omega'}\omega-\sin^2\theta\right).
\end{equation*}

相空间从 Lorentz 不变二体测度开始。末态电子能量记为
\begin{equation*}
 E_e'=\sqrt{m_{\mathrm e}^2c^4+c^2|\mathbf p'|^2}.
\end{equation*}
用空间 $\delta$ 函数消去 $\mathbf p'$ 后，
\begin{equation*}
 \dd\mathfrak P_2
 =\frac{\hbar^2\omega'}{16\pi^2E_e'}
 \dd\omega'\dd\Omega\,
 \delta\!\left(m_{\mathrm e} c^2+\hbar\omega-E_e'-\hbar\omega'\right).
\end{equation*}
令括号内的能量守恒函数为 $f(\omega')$。由
\begin{equation*}
 \mathbf p'=\frac{\hbar}{c}
 (\omega\hat{\mathbf k}-\omega'\hat{\mathbf k}')
\end{equation*}
可得物理解处
\begin{align*}
 |f'(\omega')|
 &=\frac{\hbar}{E_e'}
 \left[m_{\mathrm e} c^2+\hbar\omega(1-\cos\theta)\right]\\
 &=\frac{\hbar m_{\mathrm e} c^2}{E_e'}\frac\omega{\omega'}.
\end{align*}
因此用 $\delta[f]=\delta(\omega'-\omega'_*)/|f'(\omega'_*)|$ 完成能量积分后，
\begin{equation*}
 \frac{\dd\mathfrak P_2}{\dd\Omega}
 =\frac{\hbar}{16\pi^2m_{\mathrm e} c^2}\frac{\omega'^2}{\omega}.
\end{equation*}
正文\eqnrefs{eq:compton-invariant-flux-r14} 给出入射通量
$\mathcal I_{e\gamma}=m_{\mathrm e}\hbar\omega$。代入国际单位制下的约化振幅的截面公式
\begin{equation*}
 \dd\sigma
 =\frac{\hbar^2\overline{|\mathcal M|^2}}{4\mathcal I_{12}}
 \dd\mathfrak P_2
\end{equation*}
得到
\begin{equation*}
 \frac{\dd\sigma_C}{\dd\Omega}
 =\frac{\hbar^2}{64\pi^2m_{\mathrm e}^2c^2}
 \left(\frac{\omega'}\omega\right)^2
 \overline{|\mathcal M_C|^2}.
\end{equation*}
最后使用 $g_{\rm em}^2=4\pi\alpha$ 与
$r_e=\alpha\hbar/(m_{\mathrm e} c)$，便得到正文\eqnrefs{eq:klein-nishina-SI}。
\end{derivation}

\begin{derivation}[从 Compton 通道到双光子湮灭]
Dirac 场展开含

\begin{equation*}
 \hat\psi(x)\supset b_s(\mathbf p)u_s(p)\ee^{-\ii p\cdot x/\hbar}
 +d_s^\dagger(\mathbf p)v_s(p)\ee^{+\ii p\cdot x/\hbar}.
\end{equation*}
入射电子外腿从正频项取得 $u(p)\ee^{-\ii p\cdot x/\hbar}$；同一条费米子线改解释为出射正电子时，收缩从负频项取得 $v(p)\ee^{+\ii p\cdot x/\hbar}$。顶角的 Fourier 积分因此对应四动量解析延拓 $p\mapsto-p$，同时把 $u$ 外线替换为相应 $v$ 外线。交叉对称性正是同一时序 Green 函数在不同质量壳边界值之间的解析延拓。

对
\begin{equation*}
 e^-(p_1)+e^+(p_2)\to\gamma(k_1)+\gamma(k_2),
\end{equation*}
两个光子顶角的两种次序给出
\begin{equation*}
 {
 \mathcal M_{\gamma\gamma}
 =-g_{\rm em}^2\bar v(p_2)\left[
 \slashed\varepsilon_2^*
 \frac{\slashed p_1-\slashed k_1+m_{\mathrm e} c}{(p_1-k_1)^2-m_{\mathrm e}^2c^2}
 \slashed\varepsilon_1^*
 +\slashed\varepsilon_1^*
 \frac{\slashed p_1-\slashed k_2+m_{\mathrm e} c}{(p_1-k_2)^2-m_{\mathrm e}^2c^2}
 \slashed\varepsilon_2^*
 \right]u(p_1).}
 \end{equation*}
\end{derivation}

\begin{derivation}[双光子产生电子对的角分布]
正文\eqnrefs{eq:BW-amplitude-r68} 给出两张电子交换图的相干和。先完成自旋--偏振求和和二体相空间，只推到正文角分布；总截面的立体角积分留给下一条推导。

定义
\begin{equation*}
 s=(k_1+k_2)^2,\qquad
 t=(k_1-p_1)^2,\qquad
 u=(k_2-p_1)^2,
\end{equation*}
并使用正文\eqnrefs{eq:BW-stu-si-r12} 的 $s+t+u=2m_{\mathrm e}^2c^2$。质心系末态速率参数为
\begin{equation*}
 \beta_{\rm pair}=\sqrt{1-\frac{4m_{\mathrm e}^2c^2}{s}},
\end{equation*}
因此
\begin{equation*}
 |\mathbf k_1|=\frac{\sqrt s}{2},\qquad
 |\mathbf p_1|=\frac{\sqrt s}{2}\beta_{\rm pair}.
\end{equation*}
若 $\theta$ 是 $\mathbf k_1$ 与 $\mathbf p_1$ 的夹角，则
\begin{align*}
 t-m_{\mathrm e}^2c^2&=-\frac{s}{2}(1-\beta_{\rm pair}\cos\theta),\\
 u-m_{\mathrm e}^2c^2&=-\frac{s}{2}(1+\beta_{\rm pair}\cos\theta).
\end{align*}

Breit--Wheeler 与 Compton 属于同一个截肢四点函数的交叉道，因此八伽马迹的 Lorentz 不变量结构已经包含在\eqnrefs{eq:compton-M2-unpolarized-dp73} 中。对该结果作交叉解析延拓后，只需重新确定交叉后的 Mandelstam 不变量以及费米子自旋和所带的额外负号。

把 Compton 过程
\begin{equation*}
 e^-(p)+\gamma(k)\to e^-(p')+\gamma(k')
\end{equation*}
交叉成
\begin{equation*}
 \gamma(k_1)+\gamma(k_2)\to e^-(p_1)+e^+(p_2)
\end{equation*}
时取
\begin{equation*}
 p\to-p_2,
 \qquad
 k\to k_1,
 \qquad
 p'\to p_1,
 \qquad
 k'\to-k_2.
\end{equation*}
令
\begin{equation*}
 T\equiv t-m_{\mathrm e}^2c^2,
 \qquad
 U\equiv u-m_{\mathrm e}^2c^2.
\end{equation*}
则 Compton 中的两个不变量
\begin{equation*}
 a=p\cdot k,
 \qquad
 b=p\cdot k'
\end{equation*}
分别变成
\begin{align*}
 a&\to-p_2\cdot k_1
 =\frac{u-m_{\mathrm e}^2c^2}{2}=\frac U2,\\
 b&\to p_2\cdot k_2
 =-\frac{t-m_{\mathrm e}^2c^2}{2}=-\frac T2.
\end{align*}
第二行使用了四动量守恒 $k_2-p_2=p_1-k_1$。

平方振幅的交叉还必须处理费米子完备关系。Compton 初态电子自旋和含
\begin{equation*}
 \sum_su_s(p)\bar u_s(p)=\slashed p+m_{\mathrm e} c.
\end{equation*}
解析延拓 $p\to-p_2$ 后，
\begin{equation*}
 \slashed p+m_{\mathrm e} c
 \longrightarrow
 -\slashed p_2+m_{\mathrm e} c
 =-\bigl(\slashed p_2-m_{\mathrm e} c\bigr),
\end{equation*}
而括号内正是末态正电子的完备关系
\begin{equation*}
 \sum_sv_s(p_2)\bar v_s(p_2)=\slashed p_2-m_{\mathrm e} c.
\end{equation*}
因此交叉后的未偏振迹相对于对 Compton 不变量函数的直接代入多出一个整体负号：
\begin{equation*}
 \overline{|\mathcal M_{\rm BW}|^2}
 =-\,\overline{|\mathcal M_C|^2}
 \bigg|_{a=U/2,\,b=-T/2}.
\end{equation*}
把\eqnrefs{eq:compton-M2-unpolarized-dp73} 逐项代入：
\begin{align*}
 \overline{|\mathcal M_{\rm BW}|^2}
 =2g_{\rm em}^4\Bigg[
 &\frac UT+\frac TU
 -4m_{\mathrm e}^2c^2\left(\frac1U+\frac1T\right)\\
 &-4m_{\mathrm e}^4c^4\left(\frac1U+\frac1T\right)^2
 \Bigg].
\end{align*}
由
\begin{equation*}
 s+t+u=2m_{\mathrm e}^2c^2
 \quad\Longrightarrow\quad
 T+U=-s,
\end{equation*}
可把它等价写成
\begin{equation*}
 \overline{|\mathcal M_{\rm BW}|^2}
 =2g_{\rm em}^4\left[
 \frac UT+\frac TU
 +\frac{4m_{\mathrm e}^2c^2s}{TU}
 -\frac{4m_{\mathrm e}^4c^4s^2}{T^2U^2}
 \right].
\end{equation*}
这已经是完全 Lorentz 不变的 Breit--Wheeler 树级平方振幅。

在质心系中
\begin{equation*}
 T=-\frac s2(1-\beta_{\rm pair}\cos\theta),
 \qquad
 U=-\frac s2(1+\beta_{\rm pair}\cos\theta),
\end{equation*}
并且
\begin{equation*}
 m_{\mathrm e}^2c^2=\frac s4(1-\beta_{\rm pair}^2).
\end{equation*}
因此
\begin{align*}
 TU
 &=\frac{s^2}{4}(1-\beta_{\rm pair}^2\cos^2\theta),\\
 T^2+U^2
 &=\frac{s^2}{2}(1+\beta_{\rm pair}^2\cos^2\theta).
\end{align*}
把这些量代入上一式并通分，分子逐项为
\begin{align*}
 &2(1+\beta_{\rm pair}^2\cos^2\theta)
 (1-\beta_{\rm pair}^2\cos^2\theta)\\
 &\quad+4(1-\beta_{\rm pair}^2)
 (1-\beta_{\rm pair}^2\cos^2\theta)
 -4(1-\beta_{\rm pair}^2)^2.
\end{align*}
展开并合并同次幂后得到
\begin{equation*}
 2\left[
 1-\beta_{\rm pair}^4\cos^4\theta
 +2\beta_{\rm pair}^2(1-\beta_{\rm pair}^2)\sin^2\theta
 \right].
\end{equation*}
所以
\begin{equation*}
 \boxed{
 \overline{|\mathcal M_{\rm BW}|^2}
 =4g_{\rm em}^4
 \frac{
 1-\beta_{\rm pair}^4\cos^4\theta
 +2\beta_{\rm pair}^2(1-\beta_{\rm pair}^2)\sin^2\theta
 }{(1-\beta_{\rm pair}^2\cos^2\theta)^2}.}
\end{equation*}
这样新出现的数学内容只有交叉映射及其符号，而重复的 Dirac 迹由已经完成的 Compton 推导承担。
二体截面母\eqnrefs{eq:2to2-cross-section-master-si-r9} 给出
\begin{equation*}
 \frac{\dd\sigma_{\rm BW}}{\dd\Omega}
 =\frac{\hbar^2\beta_{\rm pair}}{64\pi^2s}
 \overline{|\mathcal M_{\rm BW}|^2}.
\end{equation*}
最后代入
\begin{equation*}
 s=\frac{4m_{\mathrm e}^2c^2}{1-\beta_{\rm pair}^2},\qquad
 g_{\rm em}^4=(4\pi\alpha)^2,\qquad
 r_e=\frac{\alpha\hbar}{m_{\mathrm e} c},
\end{equation*}
便得到正文\eqnrefs{eq:BW-dsigma-si-dp49}。
\end{derivation}

\begin{derivation}[真空极化张量的横向性]
记约化自由电子传播子为
\begin{equation*}
 S_0(\ell)=\frac{\slashed\ell+m_{\mathrm e} c}{\ell^2-m_{\mathrm e}^2c^2+\ii0^+},
 \qquad
 S_{0}^{-1}(\ell)=\slashed\ell-m_{\mathrm e} c .
\end{equation*}
由
\begin{equation*}
 \slashed q=S_0^{-1}(\ell+q)-S_0^{-1}(\ell)
\end{equation*}
收缩正文\eqnrefs{eq:vacuum-polarization-loop} 的外部动量，并把两个传播子的显式 $\ii$ 因子保留到底。利用迹的循环性得
\begin{align*}
 \ii q_\mu\Pi^{\mu\nu}(q)
 &=-g_{\rm em}^2
 \int_\ell^{(d)}
 \Tr\!\left[
 S_0(\ell+q)\slashed q\,S_0(\ell)\gamma^\nu
 \right]\\
 &=-g_{\rm em}^2
 \int_\ell^{(d)}
 \Tr\!\left[
 \bigl(S_0(\ell)-S_0(\ell+q)\bigr)\gamma^\nu
 \right].
\end{align*}
维数正则化保持圈动量平移不变；第二项作 \(\ell\mapsto\ell-q\) 后与第一项相消，因此
\begin{equation*}
 q_\mu\Pi^{\mu\nu}(q)=0.
\end{equation*}
真空 Lorentz 协变性只允许
\begin{equation*}
 \Pi^{\mu\nu}(q)=A(q^2)\eta^{\mu\nu}+B(q^2)q^\mu q^\nu .
\end{equation*}
与横向条件联立给 \(A=-q^2B\)，故得到正文\eqnrefs{eq:Pi-transverse-r12}。
\end{derivation}

\begin{derivation}[一圈真空极化的迹、积分与减法]
要从一圈张量积分得到重整化后的标量响应，先计算 Dirac 迹，再在同一维数正则化下用分部积分整理横向结构，最后施加 $\Pi_R(0)=0$ 的电荷重整化条件。保持这一步骤顺序，才能同时追踪横向性与减法常数。

记
\[
 A=\ell^2-m_{\mathrm e}^2c^2+\ii0^+,
 \qquad
 B=(\ell+q)^2-m_{\mathrm e}^2c^2+\ii0^+ .
\]
两分母参数化为
\[
 \frac1{AB}=\int_0^1\dd x\,
 \frac1{[xA+(1-x)B]^2}.
\]
令
\[
 L\equiv \ell+(1-x)q,
 \qquad
 \Delta_\Pi\equiv m_{\mathrm e}^2c^2-x(1-x)q^2,
\]
则
\[
 xA+(1-x)B=L^2-\Delta_\Pi+\ii0^+.
\]

闭合电子圈分子中的 Dirac 迹先按四伽马迹恒等式计算：
\begin{align*}
N^{\mu\nu}
&\equiv
\Tr\!\left[
\gamma^\mu(\slashed\ell+m_{\mathrm e} c)
\gamma^\nu(\slashed\ell+\slashed q+m_{\mathrm e} c)
\right]\\
&=4\left[
\ell^\mu(\ell+q)^\nu
+\ell^\nu(\ell+q)^\mu
-\eta^{\mu\nu}
\bigl(\ell\!\cdot(\ell+q)-m_{\mathrm e}^2c^2\bigr)
\right].
\end{align*}
代入
$\ell=L-(1-x)q$、$\ell+q=L+xq$。关于 $L$ 的奇次项在平移不变的维数正则化积分中为零，因此有效偶次分子为
\begin{align*}
N_{\rm even}^{\mu\nu}
=4\Bigl[
&2L^\mu L^\nu-\eta^{\mu\nu}L^2
+\eta^{\mu\nu}m_{\mathrm e}^2c^2\\
&+x(1-x)
\bigl(q^2\eta^{\mu\nu}-2q^\mu q^\nu\bigr)
\Bigr].
\end{align*}

对称积分将 $L^\mu L^\nu$ 化成与 $\eta^{\mu\nu}$ 成比例的项；要使这些项与分子中的质量项组合成横向张量，还需使用同一调节下的分部积分恒等式
\begin{align*}
0
&=\int_L^{(d)}
\frac{\partial}{\partial L_\mu}
\left[
\frac{L^\nu}{L^2-\Delta_\Pi+\ii0^+}
\right]\\
&=\eta^{\mu\nu}
\int_L^{(d)}\frac1{L^2-\Delta_\Pi+\ii0^+}
-2\int_L^{(d)}
\frac{L^\mu L^\nu}
{(L^2-\Delta_\Pi+\ii0^+)^2}.
\end{align*}
所以
\[
2\int_L^{(d)}
\frac{L^\mu L^\nu}{(L^2-\Delta_\Pi+\ii0^+)^2}
=\eta^{\mu\nu}
\int_L^{(d)}\frac1{L^2-\Delta_\Pi+\ii0^+}.
\]
另一方面
\[
\frac{L^2}{(L^2-\Delta_\Pi)^2}
=\frac1{L^2-\Delta_\Pi}
+\frac{\Delta_\Pi}{(L^2-\Delta_\Pi)^2}.
\]
因此前三项的积分恰好化成
\begin{align*}
&\int_L^{(d)}
\frac{2L^\mu L^\nu-\eta^{\mu\nu}L^2
+\eta^{\mu\nu}m_{\mathrm e}^2c^2}
{(L^2-\Delta_\Pi+\ii0^+)^2}\\
&\qquad
=\eta^{\mu\nu}
(m_{\mathrm e}^2c^2-\Delta_\Pi)
I_2(\Delta_\Pi)\\
&\qquad
=x(1-x)q^2\eta^{\mu\nu}I_2(\Delta_\Pi),
\end{align*}
其中
\[
I_2(\Delta)
\equiv
\int_L^{(d)}\frac1{(L^2-\Delta+\ii0^+)^2}.
\]
再与显式的 $q^\mu q^\nu$ 项合并，整个分子积分成为
\begin{align*}
\int_L^{(d)}
\frac{N_{\rm even}^{\mu\nu}}
{(L^2-\Delta_\Pi+\ii0^+)^2}
&=-8x(1-x)
(q^\mu q^\nu-q^2\eta^{\mu\nu})
I_2(\Delta_\Pi).
\end{align*}
横向张量因此不是事后猜出的结构，而是在同一个维数正则化积分中由分部积分恒等式直接产生。

把正文\eqnrefs{eq:vacuum-polarization-loop} 的整体因子恢复，并使用 $g_{\rm em}^2=4\pi\alpha$ 与正文\eqnrefs{eq:I2-expanded-r12}，即可从
\[
\Pi^{\mu\nu}(q)
=(q^\mu q^\nu-q^2\eta^{\mu\nu})\Pi(q^2)
\]
读出
\begin{align*}
\Pi(q^2)
=\frac{2\alpha}{\pi}
\int_0^1\dd x\,x(1-x)
\left[
\frac1\epsilon-\gamma_E+\ln4\pi
+\ln\frac{\mu_{\rm DR}^2}
{\Delta_\Pi(x,q^2)-\ii0^+}
\right]
+\order(\epsilon).
\end{align*}
在 $q^2=0$ 时，$\Delta_\Pi(x,0)=m_{\mathrm e}^2c^2$，所以同一表达式给
\begin{align*}
\Pi(0)
=\frac{2\alpha}{\pi}
\int_0^1\dd x\,x(1-x)
\left[
\frac1\epsilon-\gamma_E+\ln4\pi
+\ln\frac{\mu_{\rm DR}^2}{m_{\mathrm e}^2c^2}
\right].
\end{align*}
两式相减时，紫外极点、$\overline{\rm MS}$ 方案常数和任意重整化尺度同时消失：
\begin{align*}
\Pi_R(q^2)
&\equiv\Pi(q^2)-\Pi(0)\\
&=\frac{2\alpha}{\pi}
\int_0^1\dd x\,x(1-x)
\ln\frac{m_{\mathrm e}^2c^2}
{m_{\mathrm e}^2c^2-q^2x(1-x)-\ii0^+},
\end{align*}
这正是正文\eqnrefs{eq:Pi-renormalized-r12}。这里 $\Pi_R(0)=0$ 不是额外代数技巧，而是在壳电荷定义：零动量传递处测得的电荷被定义为 $e$，所有真空极化修正从有限的动量依赖部分开始。
\end{derivation}

\begin{derivation}[真空极化插入与有效电荷]
守恒外部电流只耦合光子传播子的横向部分。提出共同的 Lorentz 投影算符与裸分母后，一次已重整化真空极化插入给传播核乘上无量纲因子。Dyson 方程写成

\begin{equation*}
 D_T(q)=D_{0,T}(q)+D_{0,T}(q)\,\Pi_T(q)\,D_T(q).
 \end{equation*}
把第二项移到左边并提取 $D_T$：
\begin{align*}
 \left[D_{0,T}^{-1}(q)-\Pi_T(q)\right]D_T(q)&=1,\\
 D_T(q)&=\left[D_{0,T}^{-1}(q)-\Pi_T(q)\right]^{-1}.
\end{align*}
用前面定义的无量纲已重整化标量 $\Pi_R$ 参数化以后，空间样 $q^2=-Q^2$ 的横向传播子相对于裸传播子多出
\begin{equation*}
 {
 D_T(-Q^2)=\frac{D_{0,T}(-Q^2)}{1+\Pi_R(-Q^2)}.}
 \end{equation*}
因为树级电流--电流振幅只含 ``两个 QED 顶角乘一条光子传播子''，可以把同一个修正等效地吸收到耦合中，定义一圈-improved 有效 fine-结构常数
\begin{equation*}
 {
 \alpha_{\rm eff}(Q)
 \equiv
 \frac{\alpha}{1+\Pi_R(-Q^2)}.}
 \end{equation*}
\eqnrefs{eq:qed-alpha-effective-dp7} 只重求和反复出现的一圈单粒子不可约真空极化核；更高圈单粒子不可约图会继续修正 $\beta$ 函数，因此该表达式是一圈精度下的运行耦合。
\end{derivation}

\begin{derivation}[真空极化虚部与电子对阈值]
当时间样 $q^2=s>4m_{\mathrm e}^2c^2$ 时，\eqnrefs{eq:Pi-Delta-r12} 在一段 Feynman-参数区间内变负。方程

\begin{equation*}
 m_{\mathrm e}^2c^2-sx(1-x)=0
\end{equation*}
等价于
\begin{equation*}
 x^2-x+\frac{m_{\mathrm e}^2c^2}{s}=0,
\end{equation*}
根为
\begin{equation*}
 x_\pm=\frac12(1\pm\beta_{\rm pair}),
 \qquad
 \beta_{\rm pair}\equiv\sqrt{1-\frac{4m_{\mathrm e}^2c^2}{s}}.
\end{equation*}
对 $x_-<x<x_+$，分母的实部为负。由主支对数
\begin{equation*}
 \ln(-A-\ii0^+)=\ln A-\ii\pi,
 \qquad A>0,
\end{equation*}
而 \eqnrefs{eq:Pi-renormalized-r12} 中的对数是分子除以该负数，所以
\begin{equation*}
 \ln\frac{m_{\mathrm e}^2c^2}{-A-\ii0^+}
 =\ln\frac{m_{\mathrm e}^2c^2}{A}+\ii\pi.
\end{equation*}
于是
\begin{align*}
 \operatorname{Im}\Pi_R(s)
 &=2\alpha\int_{x_-}^{x_+}\dd x\,x(1-x),\\
 &=2\alpha\left[
 \frac{x^2}{2}-\frac{x^3}{3}
 \right]_{x_-}^{x_+}.
\end{align*}
利用 $x_++x_-=1$、$x_+-x_-=\beta_{\rm pair}$ 与 $x_+x_-=m_{\mathrm e}^2c^2/s$，可化为
\begin{equation*}
 {
 \operatorname{Im}\Pi_R(s)
 =\frac{\alpha}{3}\,
 \beta_{\rm pair}\left(1+\frac{2m_{\mathrm e}^2c^2}{s}\right).}
 \end{equation*}
\end{derivation}

\begin{derivation}[电子自能的结构与紫外项]
真空中除 $p^\mu$ 外没有其他优先四矢量。QED 保持宇称，自能是自旋空间中的 $4\times4$ 矩阵；Lorentz 协变性与宇称因此排除独立的 $\gamma^5$、$\gamma^\mu\gamma^5$ 与 $\sigma^{\mu\nu}$ 项，一般形式只能是

\begin{equation*}
 {
 \Sigma(\slashed p)
 =A(p^2)\slashed p+B(p^2)m_{\mathrm e} c.}
 \end{equation*}
Lorentz 协变性与 Ward 恒等式先限制顶角的允许结构，一圈积分再显式给出相应形状因子。Feynman 规范中，约化图规则给
\begin{equation*}
 \Sigma(\slashed p)
 =-\ii g_{\rm em}^2
 \int_{\ell}^{(d)}
 \frac{\gamma^\mu(\slashed p-\slashed\ell+m_{\mathrm e} c)\gamma_\mu}
 {[\ell^2+\ii0^+][(p-\ell)^2-m_{\mathrm e}^2c^2+\ii0^+]}.
 \end{equation*}
用相应的 $d$ 维 Clifford 恒等式
\begin{equation*}
 \gamma^\mu\gamma^\alpha\gamma_\mu=(2-d)\gamma^\alpha,
 \qquad
 \gamma^\mu\gamma_\mu=d\id(),
\end{equation*}
得到
\begin{equation*}
 \gamma^\mu(\slashed p-\slashed\ell+m_{\mathrm e} c)\gamma_\mu
 =(2-d)(\slashed p-\slashed\ell)+dm_{\mathrm e} c.
 \end{equation*}
用 \eqnrefs{eq:feynman-parameter-proof-r9} 合并两个分母：
\begin{align*}
 &x[(p-\ell)^2-m_{\mathrm e}^2c^2]+(1-x)\ell^2\\
 &\qquad=(\ell-xp)^2-\Delta_\Sigma(x,p^2),
\end{align*}
其中
\begin{equation*}
 {
 \Delta_\Sigma(x,p^2)
 =xm_{\mathrm e}^2c^2-x(1-x)p^2.}
 \end{equation*}
令 $L=\ell-xp$，则
\begin{equation*}
 \slashed p-\slashed\ell=(1-x)\slashed p-\slashed L.
\end{equation*}
关于 $L$ 为奇函数的 $\slashed L$ 项在对称积分中为零，因此
\begin{align*}
 \Sigma(\slashed p)
 =-\ii g_{\rm em}^2
 \int_0^1\dd x\,
 \bigl[(2-d)(1-x)\slashed p+dm_{\mathrm e} c\bigr]
 I_2(\Delta_\Sigma).
 \end{align*}
取 $d=4-2\epsilon$，只保留 \eqnrefs{eq:I2-expanded-r12} 的紫外极点：
\begin{align*}
 \Sigma_{\rm UV}(\slashed p)
 &=\frac{g_{\rm em}^2}{(4\pi)^2}\frac1\epsilon
 \int_0^1\dd x
 \left[-2(1-x)\slashed p+4m_{\mathrm e} c\right],\\
 &=\frac{\alpha}{4\pi}\frac1\epsilon
 \left(-\slashed p+4m_{\mathrm e} c\right).
 \end{align*}
因此电子场强反项与质量反项必须分别吸收 $\slashed p$ 和 $m_{\mathrm e} c$ 两种独立结构；Lorentz 协变的二点函数中没有第三个独立反项结构可供调节。
\end{derivation}

\begin{derivation}[由产生谱重建空间样真空极化]
已经由 Feynman-参数积分得到

\begin{equation*}
 \operatorname{Im}\Pi_R(s)
 =\frac{\alpha}{3}
 \beta_{\rm pair}(s)\left(1+\frac{2m_{\mathrm e}^2c^2}{s}\right),
 \qquad
 \beta_{\rm pair}(s)\equiv\sqrt{1-\frac{4m_{\mathrm e}^2c^2}{s}},
 \end{equation*}
其支撑从 $s_0=4m_{\mathrm e}^2c^2$ 开始。该阈值由光子单粒子不可约二点函数的 Cutkosky 切割决定：费米子与反费米子两条内部线同时满足
\begin{equation*}
 p_-^2=p_+^2=m_{\mathrm e}^2c^2,
 \qquad
 p_-^0>0,\quad p_+^0>0,
 \qquad
 q=p_-+p_+,
\end{equation*}
因此只有虚光子的 Lorentz 不变质量达到 $2m_{\mathrm e} c$ 后，真实 $e^-e^+$ 在壳相空间才打开。

在壳电荷重整化已固定 $\Pi_R(0)=0$。因此直接使用通用一次减法公式 \eqnrefs{eq:once-subtracted-single-cut-dp9}：
\begin{equation*}
 {
 \Pi_R(q^2)
 =\frac{q^2}{\pi}
 \int_{4m_{\mathrm e}^2c^2}^{\infty}\dd s\,
 \frac{\operatorname{Im}\Pi_R(s)}
 {s\,[s-q^2-\ii0^+]}.}
 \end{equation*}
这里减法常数已由物理重整化条件固定；割线上的谱权重则由幺正性决定，两者承担不同信息。

对空间样 $q^2=-Q^2$，$Q^2>0$，
\begin{equation*}
 \Pi_R(-Q^2)
 =-\frac{Q^2}{\pi}
 \int_{4m_{\mathrm e}^2c^2}^{\infty}\dd s\,
 \frac{\operatorname{Im}\Pi_R(s)}{s(s+Q^2)}<0,
 \end{equation*}
这与直接由 Feynman 参数得到的 \eqnrefs{eq:Pi-spacelike-r12} 符号一致。
\end{derivation}

\begin{derivation}[低动量真空极化的谱积分]
令 $a\equiv m_{\mathrm e}^2c^2$ 只在本段积分中作紧凑记号。由 \eqnrefs{eq:qed-vacpol-spacelike-dispersion-dp9}，当 $Q^2\ll a$，

\begin{align*}
 \Pi_R(-Q^2)
 &=-\frac{Q^2}{\pi}
 \int_{4a}^{\infty}\dd s\,
 \frac{\operatorname{Im}\Pi_R(s)}{s^2}
 +\order\!\left(\frac{Q^4}{a^2}\right).
 \end{align*}
把 \eqnrefs{eq:qed-vacpol-spectral-weight-dp9} 代入，并作变量
\begin{equation*}
 y\equiv\sqrt{1-\frac{4a}{s}},
 \qquad
 s=\frac{4a}{1-y^2},
 \qquad
 \dd s=\frac{8ay}{(1-y^2)^2}\dd y,
\end{equation*}
其中 $s:4a\to\infty$ 对应 $y:0\to1$。于是
\begin{align*}
 \int_{4a}^{\infty}\dd s\,
 \frac{\operatorname{Im}\Pi_R(s)}{s^2}
 &=\frac{\alpha}{3}
 \int_0^1\dd y\,
 \frac{y^2(3-y^2)}{4a},\\
 &=\frac{\alpha}{12a}
 \left[y^3-\frac{y^5}{5}\right]_0^1,\\
 &=\frac{\alpha}{15a}.
\end{align*}
因此
\begin{equation*}
 {
 \Pi_R(-Q^2)
 =-\frac{\alpha}{15\pi}
 \frac{Q^2}{m_{\mathrm e}^2c^2}
 +\order\!\left(\frac{Q^4}{m_{\mathrm e}^4c^4}\right),}
 \end{equation*}
精确恢复直接圈积分的 \eqnrefs{eq:qed-pi-smallQ-dp7}。这是一项非平凡交叉检查：一个系数由欧几里得/空间样 Feynman 参数展开得到，另一个只由时间样在壳对-产生谱权重得到。
\end{derivation}

\begin{derivation}[在壳电子为何没有真空衰变宽度]
一圈电子自能 \eqnrefs{eq:selfenergy-start-r12} 的两个内部分母对应一条光子线与一条电子线。若先以 $m_\gamma>0$ 作为光子质量的红外调节参数，两体切割的运动学阈值为

\begin{equation*}
 p^2\ge (m_{\mathrm e}+m_\gamma)^2c^2.
 \end{equation*}
因此调节参数存在时，物理电子极点 $p^2=m_{\mathrm e}^2c^2$ 与连续谱阈值分离，精确传播子具有通常的孤立极点；对当前场变量，其 LSZ 极点权重记为 $Z_{\rm pole}^{(e)}$，可由极点邻域的逆传播子导数求出。若随后采用在壳重整化，则场强反项中的 $Z_2$ 被选择来把重整化传播子的这一物理极点权重归一化为 $1$。

现在令 $m_\gamma\to0$。阈值向 $p^2=m_{\mathrm e}^2c^2$ 靠近，但这不意味着自由电子可以发出一个有限能量实光子。若假设
\begin{equation*}
 p\to p'+k,
 \qquad
 p^2=p'^2=m_{\mathrm e}^2c^2,
 \qquad
 k^2=0,
\end{equation*}
四动量守恒 $p=p'+k$ 两边平方：
\begin{align*}
 p^2
 &=(p'+k)^2,\\
 m_{\mathrm e}^2c^2
 &=m_{\mathrm e}^2c^2+2p'\!\cdot k,\\
 p'\!\cdot k&=0.
\end{align*}
对未来指向的时间样 $p'$ 与未来指向的类光 $k$，若 $k\neq0$ 则 $p'\!\cdot k>0$；所以唯一解是
\begin{equation*}
 {k^\mu=0.}
 \end{equation*}
也就是说，无质量-光子极限中自能分支点与电子质量壳相接，但没有有限相空间的 $e^-\to e^-\gamma$ 衰变通道。因此物理电子不获得通常 unstable-粒子意义下的有限衰变宽度。
\end{derivation}

\begin{derivation}[一圈顶角的参数化与动量平移]
由\eqnrefs{eq:feynman-three-denominator-r14}，分母线性组合为
\begin{align*}
xD'_e+yD_e+zD_\gamma
&=\ell^2-2\ell\cdot(xp'+yp)
+x(p'^2-m_{\mathrm e}^2c^2)+y(p^2-m_{\mathrm e}^2c^2)+\ii0^+\\
&=\ell^2-2\ell\cdot(xp'+yp)+\ii0^+.
\end{align*}
令 $L=\ell-xp'-yp$，则
\begin{equation*}
xD'_e+yD_e+zD_\gamma
=L^2-(xp'+yp)^2+\ii0^+.
\end{equation*}
由 $q=p'-p$ 与 $p^2=p'^2=m_{\mathrm e}^2c^2$，
\begin{align*}
2p\!\cdot p'&=2m_{\mathrm e}^2c^2-q^2,\\
(xp'+yp)^2
&=(x+y)^2m_{\mathrm e}^2c^2-xyq^2\\
&=(1-z)^2m_{\mathrm e}^2c^2-xyq^2
=\Delta_V(x,y,z;q^2).
\end{align*}
把这一结果代回三分母母公式即得\eqnrefs{eq:vertex-parameterized-mother-dp75}。
\end{derivation}

\begin{derivation}[顶角分子中的 Pauli 结构]

令
\begin{equation*}
 R'^\mu\equiv p'^\mu-\ell^\mu,
 \qquad
 R^\mu\equiv p^\mu-\ell^\mu.
\end{equation*}
用正文\eqnrefs{eq:vertex-shift-delta-dp75} 的 $L^\mu=\ell^\mu-xp'^\mu-yp^\mu$ 与 $x+y+z=1$，有两组等价写法
\begin{align*}
 R'
 &=-L+zp'+yq
 =-L+zp+(1-x)q,\\
 R
 &=-L+zp-xq
 =-L+zp'-(1-y)q.
 \end{align*}
所有关于 $L$ 的奇函数在平移不变的正则化中积分为零。含两个 $L$ 的项经过张量约化
$L^\rho L^\sigma\to\eta^{\rho\sigma}L^2/d$ 后只产生 $\gamma^\mu$ 结构，因此只影响 $F_1$。Pauli 系数本身紫外有限；完成这一投影后不存在可与 $d-4$ 相乘留下有限余项的紫外极点，所以随后 Pauli 分子的 Clifford 化简可在四维进行。Pauli 结构可从 $L$ 无关分子中单独抽出。

令
\begin{equation*}
 A^\mu\equiv zp^\mu+(1-x)q^\mu,
 \qquad
 B^\mu\equiv zp'^\mu-(1-y)q^\mu,
\end{equation*}
则相关矩阵是
\begin{equation*}
 N_0^\mu
 =\gamma^\alpha(\slashed A+m_{\mathrm e} c)
 \gamma^\mu(\slashed B+m_{\mathrm e} c)\gamma_\alpha.
\end{equation*}
由 Clifford 代数得到所需的五伽马矩阵夹心恒等式。由
\begin{align*}
 \gamma^\alpha\slashed A\gamma^\mu\slashed B\gamma_\alpha
 &=2\gamma^\mu\slashed B\slashed A
 -\slashed A(\gamma^\alpha\gamma^\mu\slashed B\gamma_\alpha),
\end{align*}
而
\begin{align*}
 \gamma^\alpha\gamma^\mu\slashed B\gamma_\alpha
 &=2\slashed B\gamma^\mu
 -\gamma^\mu(\gamma^\alpha\slashed B\gamma_\alpha),\\
 &=2\slashed B\gamma^\mu+2\gamma^\mu\slashed B,\\
 &=4B^\mu\id(),
\end{align*}
所以
\begin{align*}
 \gamma^\alpha\slashed A\gamma^\mu\slashed B\gamma_\alpha
 &=2\gamma^\mu\slashed B\slashed A-4B^\mu\slashed A,\\
 &=-2\slashed B\gamma^\mu\slashed A.
 \end{align*}
最后一步使用
$\slashed B\gamma^\mu=2B^\mu-\gamma^\mu\slashed B$。其余三个收缩分别计算为
\begin{align*}
 \gamma^\alpha\gamma^\mu\slashed B\gamma_\alpha
 &=\gamma^\alpha\gamma^\mu\gamma^\beta\gamma_\alpha B_\beta
 =4\eta^{\mu\beta}B_\beta
 =4B^\mu,\\
 \gamma^\alpha\slashed A\gamma^\mu\gamma_\alpha
 &=A_\beta\gamma^\alpha\gamma^\beta\gamma^\mu\gamma_\alpha
 =4A^\mu,\\
 \gamma^\alpha\gamma^\mu\gamma_\alpha
 &=(2\eta^{\alpha\mu}-\gamma^\mu\gamma^\alpha)\gamma_\alpha
 =2\gamma^\mu-4\gamma^\mu
 =-2\gamma^\mu.
\end{align*}
于是
\begin{equation*}
 N_0^\mu
 =-2\slashed B\gamma^\mu\slashed A
 +4m_{\mathrm e} c(A+B)^\mu
 -2m_{\mathrm e}^2c^2\gamma^\mu.
 \end{equation*}

把它夹在 $\bar u(p')$ 与 $u(p)$ 之间，并把可用的在壳 Dirac 方程移到最外侧，使用
\begin{equation*}
 \slashed B=z\slashed p'-(1-y)\slashed q,
 \qquad
 \slashed A=z\slashed p+(1-x)\slashed q.
\end{equation*}
因此
\begin{align*}
 \bar u\slashed B\gamma^\mu\slashed A u
 =&\,\bar u\Bigl[
 z^2m_{\mathrm e}^2c^2\gamma^\mu
 +zm_{\mathrm e} c(1-x)\gamma^\mu\slashed q\\
 &\qquad-zm_{\mathrm e} c(1-y)\slashed q\gamma^\mu
 -(1-x)(1-y)\slashed q\gamma^\mu\slashed q
 \Bigr]u.
\end{align*}
利用
\begin{align*}
 \gamma^\mu\slashed q
 &=q^\mu-\ii\sigma^{\mu\nu}q_\nu,\\
 \slashed q\gamma^\mu
 &=q^\mu+\ii\sigma^{\mu\nu}q_\nu,
\end{align*}
线性 $q$ 的两项给
\begin{equation*}
 zm_{\mathrm e} c\left[(y-x)q^\mu
 -\ii(2-x-y)\sigma^{\mu\nu}q_\nu\right]
 =zm_{\mathrm e} c\left[(y-x)q^\mu
 -\ii(1+z)\sigma^{\mu\nu}q_\nu\right].
\end{equation*}
而
\begin{equation*}
 \slashed q\gamma^\mu\slashed q
 =2q^\mu\slashed q-q^2\gamma^\mu
\end{equation*}
在在壳旋量夹心中满足
\begin{align*}
 \bar u(p')\slashed q\,u(p)
 &=\bar u(p')(\slashed p'-\slashed p)u(p),\\
 &=m_{\mathrm e} c\,\bar uu-m_{\mathrm e} c\,\bar uu=0,
\end{align*}
故该项只进入 $\gamma^\mu$，不贡献 Pauli 结构。

再处理 $N_0^\mu$ 的第二项：
\begin{equation*}
 A+B=z(p'+p)+(y-x)q.
\end{equation*}
其中 $(y-x)q^\mu$ 在对称 simplex 积分 $x\leftrightarrow y$ 下为奇函数，积分为零。对 $(p'+p)^\mu$ 使用 Gordon 恒等\eqnrefs{eq:gordon-identity-r14}，
\begin{equation*}
 (p'+p)^\mu
 \doteq 2m_{\mathrm e} c\gamma^\mu-\ii\sigma^{\mu\nu}q_\nu,
\end{equation*}
其中 $\doteq$ 表示两侧夹在同一对在壳旋量之间时相等。因此，第一项
$-2\slashed B\gamma^\mu\slashed A$ 的 Pauli 部分是
\begin{equation*}
 +2\ii zm_{\mathrm e} c(1+z)\sigma^{\mu\nu}q_\nu,
\end{equation*}
第二项 $4m_{\mathrm e} c(A+B)^\mu$ 的 Pauli 部分是
\begin{equation*}
 -4\ii zm_{\mathrm e} c\sigma^{\mu\nu}q_\nu.
\end{equation*}
相加得到
\begin{equation*}
 N_0^\mu\Big|_{\rm Pauli}
 =-2\ii z(1-z)m_{\mathrm e} c\,\sigma^{\mu\nu}q_\nu.
\end{equation*}
按 \eqnrefs{eq:form-factor-definition-r14} 的标准基底写成
\begin{equation*}
 N_0^\mu\Big|_{\rm Pauli}
 =-N_2\frac{\ii\sigma^{\mu\nu}q_\nu}{2m_{\mathrm e} c},
\end{equation*}
便得到
\begin{equation*}
 {
 N_2(x,y,z)=4z(1-z)m_{\mathrm e}^2c^2.}
 \end{equation*}
该结果不依赖圈动量，因此 $F_2$ 的动量积分比 $F_1$ 更收敛。
\end{derivation}

\begin{derivation}[零动量 Pauli 形状因子的积分]

由参数化后的顶角积分和刚得到的 $N_2=4z(1-z)m_{\mathrm e}^2c^2$，一圈 Pauli 形状因子是
\begin{equation*}
 F_2^{\rm 1L}(q^2)
 =2\ii g_{\rm em}^2
 \int_0^1\dd x\dd y\dd z\,
 \delta(x+y+z-1)
 \int\frac{\dd^4 L}{(2\pi)^4}
 \frac{N_2}{(L^2-\Delta_V+\ii0^+)^3}.
 \end{equation*}
Pauli 投影对应的积分已经紫外有限；完成正则化代数后可以取 $d=4$。

对其中的标量圈积分直接作 Wick 转动 $L^0=\ii L_E^0$。由于 $\dd^4 L=\ii\dd^4 L_E$ 且 $L^2-\Delta=-(L_E^2+\Delta)$，有
\begin{equation*}
 \int\frac{\dd^4 L}{(2\pi)^4}\frac1{(L^2-\Delta+\ii0^+)^3}
 =-\ii\int\frac{\dd^4 L_E}{(2\pi)^4}\frac1{(L_E^2+\Delta)^3}.
\end{equation*}
按正文的 $\Omega_d\equiv\int\dd\Omega_{d-1}$ 记号，四维 Euclid 空间的单位三球面面积为 $\Omega_4=2\pi^2$；写 $R=|L_E|$，
\begin{align*}
 \int\frac{\dd^4 L}{(2\pi)^4}\frac1{(L^2-\Delta+\ii0^+)^3}
 &=-\ii\frac{2\pi^2}{(2\pi)^4}
 \int_0^\infty\dd R\,
 \frac{R^3}{(R^2+\Delta)^3},\\
 &=-\frac{\ii}{8\pi^2}\frac12
 \int_0^\infty\dd u\,
 \frac{u}{(u+\Delta)^3},
 \qquad u\equiv R^2.
\end{align*}
再令 $u=\Delta v$：
\begin{align*}
 \int_0^\infty\dd u\frac{u}{(u+\Delta)^3}
 &=\frac1\Delta
 \int_0^\infty\dd v\frac{v}{(1+v)^3},\\
 &=\frac1\Delta
 \left[\frac12\right],
\end{align*}
其中最后一项可由 $w=1/(1+v)$ 直接积分得到。因此
\begin{equation*}
 {
 \int\frac{\dd^4 L}{(2\pi)^4}\frac1{(L^2-\Delta+\ii0^+)^3}
 =-\frac{\ii}{32\pi^2}\frac1\Delta.}
 \end{equation*}

现在取 $q^2=0$。由参数化分母
\begin{equation*}
 \Delta_V=(1-z)^2m_{\mathrm e}^2c^2,
\end{equation*}
所以
\begin{equation*}
 \frac{N_2}{\Delta_V}
 =\frac{4z(1-z)m_{\mathrm e}^2c^2}{(1-z)^2m_{\mathrm e}^2c^2}
 =\frac{4z}{1-z}.
\end{equation*}
把上面求得的标量圈积分代入 Pauli 形状因子的母积分：
\begin{align*}
 F_2^{\rm 1L}(0)
 &=\frac{g_{\rm em}^2}{16\pi^2}
 \int_0^1\dd x\dd y\dd z\,
 \delta(x+y+z-1)
 \frac{4z}{1-z}.
\end{align*}
固定 $z$ 后，单纯形上 $x+y=1-z$，故
\begin{align*}
 \int_0^1\dd x\dd y\,
 \delta(x+y+z-1)
 &=\int_0^{1-z}\dd x=1-z.
\end{align*}
该几何因子恰好消去被积函数中的 $1-z$：
\begin{align*}
 F_2^{\rm 1L}(0)
 &=\frac{g_{\rm em}^2}{16\pi^2}
 \int_0^1\dd z\,4z,\\
 &=\frac{g_{\rm em}^2}{8\pi^2},\\
 &={\frac{\alpha}{2\pi}},
 \end{align*}
最后一步用了 $g_{\rm em}^2=4\pi\alpha$。整个计算中 $m_{\mathrm e}^2c^2$ 在分子与分母中显式相消，因此 $F_2(0)$ 自然保持无量纲。
\end{derivation}

\begin{derivation}[Pauli 形状因子与反常磁矩]
最小耦合 Dirac Hamilton 量的低能展开给出自旋磁矩项

\begin{equation*}
 H_{B}^{\rm Dirac}
 =-\frac{q_{\mathrm e}\hbar}{2m_{\mathrm e}}\boldsymbol\sigma_{\rm P}\cdot\mathbf B
 =-\frac{q_{\mathrm e}}{m_{\mathrm e}}\mathbf S_{1/2}\cdot\mathbf B,
\end{equation*}
其中 $\mathbf S_{1/2}=\hbar\boldsymbol\sigma_{\rm P}/2$。正文已经用
$H_B=-(g_{\rm mag}q_{\mathrm e}/2m_{\mathrm e})\mathbf S_{1/2}\cdot\mathbf B$ 固定了旋磁因子的物理含义；这里仅比较顶角中两个磁张量的系数。

由 Gordon 恒等式，$F_1\gamma^\mu$ 自身已经含有
$F_1\,\ii\sigma^{\mu\nu}q_\nu/(2m_{\mathrm e} c)$ 磁张量；Pauli 形状因子又增加
$F_2\,\ii\sigma^{\mu\nu}q_\nu/(2m_{\mathrm e} c)$。因此静态极限 $q^2\to0$ 的磁耦合相对于树级 Dirac 结果乘上
$F_1(0)+F_2(0)$：
\begin{equation*}
 {
 g_{\rm mag}=2\,[F_1(0)+F_2(0)].}
 \end{equation*}
物理电荷归一化与 Ward--Takahashi 恒等式共同固定
\begin{equation*}
 F_1(0)=1.
\end{equation*}
正文\eqnrefs{eq:ae-form-factor-r68} 已把相对 Dirac 值的偏移记为 $a_\ee^{\rm anom}=(g_{\rm mag}-2)/2$，因此
\begin{equation*}
 a_\ee^{\rm anom}=F_2(0).
\end{equation*}
结合 \eqnrefs{eq:schwinger-F2-r14} 得到 Schwinger 领头修正
\begin{equation*}
 {
 a_\ee^{\rm anom}
 =\frac{\alpha}{2\pi}+\order(\alpha^2),
 \qquad
 g_{\rm mag}
 =2+\frac{\alpha}{\pi}+\order(\alpha^2).}
 \end{equation*}
\end{derivation}

\begin{derivation}[空间样动量下的 Pauli 形状因子]

同一计算还可以保留 $q^2\neq0$。正文已经取空间样转移动量 $Q^2=-q^2>0$；在国际单位制中 $Q$ 的量纲是动量。代入上面已经直接算出的标量圈积分，得到
\begin{equation*}
 F_2^{\rm 1L}(q^2)
 =\frac{g_{\rm em}^2}{16\pi^2}
 \int_0^1\dd x\dd y\dd z\,
 \delta(x+y+z-1)
 \frac{4z(1-z)m_{\mathrm e}^2c^2}
 {(1-z)^2m_{\mathrm e}^2c^2-xyq^2}.
 \end{equation*}
令
\begin{equation*}
 w\equiv1-z=x+y,
 \qquad
 u_F\equiv\frac{x}{x+y},
 \qquad
 x=w u_F,
 \qquad
 y=w(1-u_F).
\end{equation*}
这里 $w,u_F\in[0,1]$ 都是无量纲变量。雅可比因子是
\begin{equation*}
 \left|\det\frac{\partial(x,y)}{\partial(w,u_F)}\right|=w.
\end{equation*}
于是分母与分子分别变成
\begin{align*}
 (1-z)^2m_{\mathrm e}^2c^2-xyq^2
 &=w^2[m_{\mathrm e}^2c^2-u_F(1-u_F)q^2],\\
 4z(1-z)m_{\mathrm e}^2c^2
 &=4(1-w)w\,m_{\mathrm e}^2c^2.
\end{align*}
连同测度的额外 $w$，所有 $w^2$ 相消：
\begin{align*}
 F_2^{\rm 1L}(q^2)
 &=\frac{g_{\rm em}^2}{16\pi^2}
 \int_0^1\dd u_F
 \frac{m_{\mathrm e}^2c^2}{m_{\mathrm e}^2c^2-u_F(1-u_F)q^2}
 \int_0^1\dd w\,4(1-w),\\
 &={
 \frac{\alpha}{2\pi}
 \int_0^1\dd u_F\,
 \frac{m_{\mathrm e}^2c^2}
 {m_{\mathrm e}^2c^2-u_F(1-u_F)q^2}}.
 \end{align*}
在空间样区域 $q^2=-Q^2$，被积函数没有极点：
\begin{equation*}
 \frac{F_2^{\rm 1L}(-Q^2)}{F_2^{\rm 1L}(0)}
 =\int_0^1\dd u_F\,
 \frac1{1+[Q/(m_{\mathrm e} c)]^2u_F(1-u_F)}.
 \end{equation*}
\end{derivation}

\begin{derivation}[真空极化怎样修正静态势能]
取点状原子核电荷为 $+Ze$、电子电荷为 $q_{\mathrm e}=-e$。静态 Coulomb 势能为

\begin{equation*}
 V_C(r)=-\frac{Ze^2}{4\pi\varepsilon_0r}
 =-\frac{Z\alpha\hbar c}{r}.
 \end{equation*}
空间 Fourier 变换写为
\begin{equation*}
 V(\mathbf r)
 =\int\frac{\dd^3 q}{(2\pi\hbar)^3}
 \ee^{+\ii\mathbf q\cdot\mathbf r/\hbar}
 \widetilde V(\mathbf q).
 \end{equation*}
由
\begin{equation*}
 \int\frac{\dd^3 q}{(2\pi\hbar)^3}
 \ee^{+\ii\mathbf q\cdot\mathbf r/\hbar}
 \frac{\hbar^2}{\mathbf q^2}
 =\frac1{4\pi r},
\end{equation*}
可知
\begin{equation*}
 \widetilde V_C(\mathbf q)
 =-\frac{Ze^2\hbar^2}{\varepsilon_0\mathbf q^2}.
 \end{equation*}

静态交换满足 $q^0=0$，因而 $q^2=-Q^2$，其中 $Q\equiv|\mathbf q|$。按真空极化的符号约定，定义空间样屏蔽函数
\begin{equation*}
 \widehat\Pi(Q^2)
 \equiv-\Pi_R(-Q^2)
 =\frac{2\alpha}{\pi}
 \int_0^1\dd x\,x(1-x)
 \ln\left[1+\frac{Q^2}{m_{\mathrm e}^2c^2}x(1-x)\right].
 \end{equation*}
Dyson-resummed 横向光子分母是 $Q^2[1+\Pi_R(-Q^2)]$。保留一圈的一阶修正，
\begin{equation*}
 \frac1{Q^2[1+\Pi_R(-Q^2)]}
 =\frac1{Q^2}[1+\widehat\Pi(Q^2)]+
 \order(\alpha^2).
\end{equation*}
因此
\begin{equation*}
 \delta\widetilde V_U(\mathbf q)
 =\widetilde V_C(\mathbf q)\widehat\Pi(Q^2).
 \end{equation*}

三维 Fourier 变换可借助一次减除的色散表示。$\Pi_R(z)/z$ 在复 $z$ 平面除 $z\ge4m_{\mathrm e}^2c^2$ 的支割外解析，并且 $\Pi_R(0)=0$。对包围支割的 Cauchy 围道，
\begin{equation*}
 \frac{\Pi_R(z)}{z}
 =\frac1{2\pi\ii}
 \int_{4m_{\mathrm e}^2c^2}^{\infty}\dd s\,
 \frac{\Pi_R(s+\ii0)-\Pi_R(s-\ii0)}{s(s-z)}.
\end{equation*}
由于不连续度为 $2\ii\operatorname{Im}\Pi_R(s)$，取 $z=-Q^2$ 得
\begin{equation*}
 \widehat\Pi(Q^2)
 =\frac{Q^2}{\pi}
 \int_{4m_{\mathrm e}^2c^2}^{\infty}\dd s\,
 \frac{\operatorname{Im}\Pi_R(s)}{s(s+Q^2)}.
 \end{equation*}
再代入 \eqnrefs{eq:ImPi-r12}：
\begin{equation*}
 \frac{\widehat\Pi(Q^2)}{Q^2}
 =\frac{\alpha}{3\pi}
 \int_{4m_{\mathrm e}^2c^2}^{\infty}\frac{\dd s}{s(s+Q^2)}
 \sqrt{1-\frac{4m_{\mathrm e}^2c^2}{s}}
 \left(1+\frac{2m_{\mathrm e}^2c^2}{s}\right).
 \end{equation*}

相应的 Fourier 变换为
\begin{equation*}
 I_s(r)
 \equiv
 \int\frac{\dd^3 q}{(2\pi\hbar)^3}
 \ee^{+\ii\mathbf q\cdot\mathbf r/\hbar}
 \frac{\hbar^2}{\mathbf q^2+s}.
\end{equation*}
取球坐标并把极轴沿 $\mathbf r$，角积分为
\begin{equation*}
 \int\dd\Omega_q\,\ee^{\ii Qr\cos\theta/\hbar}
 =4\pi\frac{\sin(Qr/\hbar)}{Qr/\hbar}.
\end{equation*}
所以
\begin{equation*}
 I_s(r)
 =\frac1{2\pi^2r}
 \int_0^\infty\dd Q\,
 \frac{Q\sin(Qr/\hbar)}{Q^2+s}.
\end{equation*}
将 $\sin$ 写成指数并在上半平面闭合积分轮廓，唯一极点在 $Q=+\ii\sqrt s$；留数定理给
\begin{equation*}
 \int_0^\infty\dd Q\,
 \frac{Q\sin(Qr/\hbar)}{Q^2+s}
 =\frac\pi2\ee^{-\sqrt s\,r/\hbar},
\end{equation*}
因此
\begin{equation*}
 {
 I_s(r)=\frac{\ee^{-\sqrt s\,r/\hbar}}{4\pi r}.}
 \end{equation*}
把 \eqnrefs{eq:Pi-dispersion-explicit-r12} 插入 \eqnrefs{eq:Uehling-momentum-correction-r12}，交换 $s$ 与 $\mathbf q$ 积分，并使用上面得到的三维 Fourier 积分：
\begin{align*}
 \delta V_U(r)
 &=V_C(r)\frac{\alpha}{3\pi}
 \int_{4m_{\mathrm e}^2c^2}^{\infty}\frac{\dd s}{s}
 \ee^{-\sqrt s\,r/\hbar}
 \sqrt{1-\frac{4m_{\mathrm e}^2c^2}{s}}
 \left(1+\frac{2m_{\mathrm e}^2c^2}{s}\right).
\end{align*}
令
\begin{equation*}
 \sqrt s=2m_{\mathrm e} c\,t,
 \qquad t\in[1,\infty),
 \qquad \frac{\dd s}{s}=\frac{2\dd t}{t},
\end{equation*}
并定义约化 Compton 波长 $\lambda_C\equiv\hbar/(m_{\mathrm e} c)$，得到经典 Uehling 修正：
\begin{equation*}
 {
 \delta V_U(r)
 =V_C(r)\frac{2\alpha}{3\pi}
 \int_1^\infty\dd t\,
 \ee^{-2rt/\lambda_C}
 \left(1+\frac1{2t^2}\right)
 \frac{\sqrt{t^2-1}}{t^2}.}
 \end{equation*}
\end{derivation}

\begin{derivation}[平面波背景中的 Volkov 解]
强经典激光背景不应在一开始按光子数截断。令经典四势只依赖平面波相位

\begin{equation*}
 \phi\equiv\kappa\cdot x,
 \qquad
 \kappa^2=0,
 \qquad
 \kappa\cdot A(\phi)=0,
 \end{equation*}
其中 $\kappa^\mu=(\omega_L/c,\mathbf k_L)$ 是四波矢，量纲为 $\mathrm{m^{-1}}$；$p^\mu$ 仍是电子物理四动量。Dirac 方程写成
\begin{equation*}
 [\ii\hbar\slashed\partial-q_{\mathrm e}\slashed A(\phi)-m_{\mathrm e} c]\psi(x)=0.
 \end{equation*}
原式整体除以 $c$ 后，三项都具有动量量纲。

取自由在壳自旋子 $u_s(p)$，满足
\begin{equation*}
 (\slashed p-m_{\mathrm e} c)u_s(p)=0.
\end{equation*}
定义
\begin{align*}
 R_p(\phi)
 &\equiv
 1+\frac{q_{\mathrm e}}{2p\cdot\kappa}
 \slashed\kappa\slashed A(\phi),\\
 F_p(\phi)
 &\equiv
 \frac{q_{\mathrm e} p\cdot A(\phi)}{p\cdot\kappa}
 -\frac{q_{\mathrm e}^2A^2(\phi)}{2p\cdot\kappa},\\
 S_p(x)
 &\equiv p\cdot x+\int^{\phi}\dd\varphi\,F_p(\varphi).
 \end{align*}
量纲检查很直接：$p\cdot\kappa$ 是动量/长度，$q_{\mathrm e} A$ 是动量，因此 $R_p$ 无量纲；$F_p$ 具有作用量量纲，因为相位变量 $\phi$ 无量纲，所以 $S_p/\hbar$ 无量纲。

候选解为
\begin{equation*}
 {
 \psi_{p,s}^{\rm V}(x)
 =R_p(\phi)u_s(p)
 \exp\left[-\frac{\ii}{\hbar}S_p(x)\right].}
 \end{equation*}
将该形式直接代回 Dirac 方程。首先
\begin{equation*}
 \partial_\mu S_p=p_\mu+\kappa_\mu F_p,
 \qquad
 \partial_\mu R_p=\kappa_\mu R_p'(\phi).
\end{equation*}
由于 $\slashed\kappa^2=\kappa^2=0$，
\begin{equation*}
 \slashed\partial R_p
 =\slashed\kappa R_p'
 =\frac{q_{\mathrm e}}{2p\cdot\kappa}
 \slashed\kappa\slashed\kappa\slashed A'(\phi)
 =0.
 \end{equation*}
所以\eqnrefs{eq:Volkov-Dirac-r12} 作用在候选解上只剩
\begin{equation*}
 [\slashed p+\slashed\kappa F_p-q_{\mathrm e}\slashed A-m_{\mathrm e} c]R_pu_s.
 \end{equation*}
需要证明它为零。先计算 $\slashed pR_pu_s$。利用
\begin{equation*}
 \slashed p\slashed\kappa=2p\cdot\kappa-\slashed\kappa\slashed p,
 \qquad
 \slashed p\slashed A=2p\cdot A-\slashed A\slashed p,
\end{equation*}
以及 $\slashed p u_s=m_{\mathrm e} cu_s$：
\begin{align*}
 \slashed pR_pu_s
 &=m_{\mathrm e} cu_s+\frac{q_{\mathrm e}}{2p\cdot\kappa}
 \slashed p\slashed\kappa\slashed A u_s,\\
 &=m_{\mathrm e} cu_s+q_{\mathrm e}\slashed A u_s
 -\frac{q_{\mathrm e}}{2p\cdot\kappa}
 \slashed\kappa\slashed p\slashed A u_s,\\
 &=m_{\mathrm e} cR_pu_s+q_{\mathrm e}\slashed A u_s
 -q_{\mathrm e}\frac{p\cdot A}{p\cdot\kappa}\slashed\kappa u_s.
 \end{align*}
另一方面，$\kappa\cdot A=0$ 给
\begin{equation*}
 \slashed A\slashed\kappa\slashed A
 =-A^2\slashed\kappa,
\end{equation*}
于是
\begin{equation*}
 q_{\mathrm e}\slashed A R_pu_s
 =q_{\mathrm e}\slashed A u_s
 -\frac{q_{\mathrm e}^2A^2}{2p\cdot\kappa}\slashed\kappa u_s.
 \end{equation*}
同时 $\slashed\kappa R_p=\slashed\kappa$，故
\begin{equation*}
 \slashed\kappa F_p R_pu_s
 =\left[
 q_{\mathrm e}\frac{p\cdot A}{p\cdot\kappa}
 -\frac{q_{\mathrm e}^2A^2}{2p\cdot\kappa}
 \right]\slashed\kappa u_s.
 \end{equation*}
$\slashed pR_pu_s$、$q_{\mathrm e}\slashed A R_pu_s$ 与 $\slashed\kappa F_pR_pu_s$ 三个结果相加后，$q_{\mathrm e}\slashed A$、$p\cdot A$ 与 $A^2$ 三类项逐项相消，最后只剩
\begin{equation*}
 m_{\mathrm e} cR_pu_s-m_{\mathrm e} cR_pu_s=0.
\end{equation*}
因此 \eqnrefs{eq:Volkov-solution-r12} 确实解含任意强度经典平面波的 Dirac 方程。

对周期平面波且选规范使 $\langle A\rangle_\phi=0$，长时间平均相位中的线性项消失。定义 Volkov 准动量
\begin{equation*}
 {
 \bar p^\mu
 =p^\mu
 -\frac{q_{\mathrm e}^2\langle A^2\rangle_\phi}{2p\cdot\kappa}\kappa^\mu.}
 \end{equation*}
因为 $\kappa^2=0$，
\begin{align*}
 \bar p^2
 &=p^2-q_{\mathrm e}^2\langle A^2\rangle_\phi,\\
 &=m_{\mathrm e}^2c^2-q_{\mathrm e}^2\langle A^2\rangle_\phi.
\end{align*}
横向平面波有 $A^2<0$，所以准动量对应的有效质量增大。若定义峰值归一化矢势
\begin{equation*}
 a_0\equiv\frac{|q_{\mathrm e}|A_0}{m_{\mathrm e} c},
\end{equation*}
则圆偏振的 $\langle A^2\rangle=-A_0^2$ 给
\begin{equation*}
 {m_*^{\rm circ}=m_{\mathrm e}\sqrt{1+a_0^2},}
\end{equation*}
而线偏振可取 $A^\mu(\phi)=\mathcal A^\mu\cos\phi$，并定义正的峰值幅度 $A_0\equiv\sqrt{-\mathcal A^2}$；于是 $\langle A^2\rangle=\mathcal A^2/2=-A_0^2/2$，故
\begin{equation*}
 {m_*^{\rm lin}=m_{\mathrm e}\sqrt{1+\frac{a_0^2}{2}}.}
 \end{equation*}
这些是完整平面波 Dirac 动力学的结果；只有在随后再取 $a_0\ll1$ 时，才可以把它展开成普通的弱场多光子微扰绘景。
\end{derivation}
